\section{Part 11a-3: the iteration schedule, the inductive lemma, initialisation and parameter selection}\label{bpP:p11a3}
\begingroup\sloppy

\def\Ddom#1#2#3{\mathcal D(#1,#2;#3)}
\def\nrmm#1#2{\lVert #1\rVert_{#2}}
\def\LamD#1{\Lambda_{#1}}
\def\dgm#1{\lvert #1\rvert_{\mathrm{dg}}}
\def\opnm#1{\lvert #1\rvert_{\mathrm{op}}}
\def\Sh{\mathcal S}
\def\cutT{\mathcal T}

This is the third of the five files of Part~11a (split plan \texttt{PLAN-11a.md}). Files~11a-1 and~11a-2 prove the two halves of one KAM step: the \emph{measure} of the parameters whose frequency violates a power-law small-divisor condition, and the \emph{analytic} step (Fourier and Taylor cutoffs, homological equation, simple canonical map) on a complex neighbourhood of the retained parameter set. This file runs the step infinitely often. It fixes the schedule of every quantity (cutoff $T_\nu$, size $\varepsilon_\nu$, strip width $r_\nu$, action radius $s_\nu$, parameter collar $t_\nu$), proves that the hypotheses of File~11a-2 and File~11a-1 hold at every stage, chooses the free parameters explicitly, and extracts the set of parameters on which the frequencies converge to a $(\gamma,\tau)$-Diophantine vector.

\paragraph{The mechanism.} R\"ussmann's iteration is only \emph{geometric}: the remainder shrinks by a fixed factor $\varpi<1$ per step (it is not superlinear), because the cutoff error and the loss of action radius contribute terms linear in $\varepsilon_\nu$ (Remark~\ref{bpP:p11a:step-reading}(d)). Geometric decay is enough because the three competing growth rates are all \emph{fixed geometric ratios}: the inverse divisor $D_\nu^{-1}=\gamma^{-1}T_\nu^{\tau}$ grows by $2^{\tau}$ per step, the inverse action radius by $1/q$, the inverse collar $t_\nu^{-1}$ by $2^{\tau+1}$; the decay ratio $\varpi$ and the shrink factor $q$ of the action radius are chosen \emph{once}, depending only on $(N,\mu,\tau)$, so small that $\varpi$ beats every one of these growth factors at every stage (Lemma~\ref{bpP:p11a:schedule-ratios}). What remains is a finite list of inequalities at stage $\nu=0$, polynomial in the data (Definition~\ref{bpP:p11a:adm-params}); this is where the \emph{shrinking-domain budget} is spent (the term $\tfrac N2C_1s_0^2$ of the initial remainder, Lemma~\ref{bpP:p11a:initial}).

\paragraph{What is proved.}
\begin{itemize}
\item \texttt{geom-recursion}, \texttt{schedule-consts}, \texttt{schedule-ratios}, \texttt{schedule}, \texttt{strip-budget}, \texttt{schedule-props} (Lemma~\ref{bpP:p11a:geom-recursion}, Definitions~\ref{bpP:p11a:schedule-consts}, \ref{bpP:p11a:schedule}, Lemmas~\ref{bpP:p11a:schedule-ratios}, \ref{bpP:p11a:strip-budget}, \ref{bpP:p11a:schedule-props}): the geometric-recursion version of track~K's iteration lemma and the schedule, with every exponent inequality of Rus01 (13.1.2), (13.1.6) replaced by an explicit one.
\item \texttt{smooth-cutoff}, \texttt{finite-extension} (Lemmas~\ref{bpP:p11a:smooth-cutoff}, \ref{bpP:p11a:finite-extension}): the chain extension $\tilde\omega_\nu$ of the stage-$\nu$ frequency to a $C^\mu$ map on the real neighbourhood of $K$ by explicit mollified cutoffs, with the closeness to $\omega$ that File~11a-1 needs (deviation D-7: it \emph{works}; Rus01 \S\S19--20 are not needed).
\item \texttt{initial}, \texttt{adm-params}, \texttt{adm-choice} (Lemma~\ref{bpP:p11a:initial}, Definition~\ref{bpP:p11a:adm-params}, Proposition~\ref{bpP:p11a:adm-choice}): the translation of Rus01 \S14, the initial remainder, the explicit admissible parameters $(T_0,t_0,s_0,\varepsilon_{\mathrm{adm}})$ and the proof that they satisfy the inequalities (A1)--(A7).
\item \texttt{finite-divisor}, \texttt{finite-measure}, \texttt{inductive-lemma} (Lemmas~\ref{bpP:p11a:finite-divisor}, \ref{bpP:p11a:finite-measure}, Proposition~\ref{bpP:p11a:inductive-lemma}): the \emph{finite-stage divisor invariant} (it replaces the limit-based tail (I3) of File~11a-1, see Remark~\ref{bpP:p11a:i3-corrected}), the nonemptiness of the retained set by the measure estimate of File~11a-1, and the inductive lemma.
\item \texttt{parameter-selection}, \texttt{exported-sums} (Theorem~\ref{bpP:p11a:parameter-selection}, Lemma~\ref{bpP:p11a:exported-sums}): the retained sets, the limit frequency, its Diophantine property and the measure loss $\le\varepsilon^*$; the displacement, derivative and Lipschitz budgets that File~11a-4 consumes.
\end{itemize}

\paragraph{Verdicts.} (R1) The admissible size $\varepsilon_{\mathrm{adm}}$ is an explicit minimum of three expressions (Remark~\ref{bpP:p11a:adm-branches}); it is \emph{not} a monomial, and its dependence on the data is controlled only when \emph{all} the data are controlled: Proposition~\ref{bpP:p11a:adm-budget} gives a lower bound $\varepsilon_{\mathrm{adm}}\ge c\,u^{B}$ along any family of data with polynomially controlled $\vartheta,s,\beta,\varepsilon^*$ from below and $M,M_0,C_1,d$ from above, with an explicit exponent $B$ (independent of the iteration index), and Corollary~\ref{bpP:p11a:adm-budget-beta} is the special case needed by Part~11 (everything fixed except the amount $\beta$). Without a bound on $M$ or on $C_1$ no such lower bound holds (review P11A3-01; see the Remark after the Proposition). Numerics: Section~\ref{bpP:p11a:sec-num}. (D-7) The chain extension by mollified cutoffs is valid under the exponent condition $g_3=\varrho\,2^{(\tau+1)\mu}\le2^{-5}$, which the choice of $\varpi$ enforces (this is the analogue of Rus01's $\tau\mu<\tau_0$, Lemma~\ref{bpP:p11a:schedule-ratios}); no fall-back to Rus01 \S\S19--20 is needed.

%%%%%%%%%%%%%%%%%%%%%%%%%%%%%%%%%%%%%%%%%%%%%%%%%%%%%%%%%%%%%%%%%%%%%%%%%%%%%%%%
\subsection{Standing data, conventions, external facts}\label{bpP:p11a:s3-conv}

All conventions are those of Files~11a-1 and~11a-2: angles of period $1$, characters $e^{2\pi ik\cdot\theta}$, max norms $\lvert\cdot\rvert_\infty$ on $\R^N,\C^N$ with induced operator norms, $\lvert k\rvert_1=\sum\lvert k_i\rvert$, domains $T_r$, $\mathcal D(r,s;B)=\{(\theta,\eta,b):\theta\in T_r,\ \lvert\eta\rvert_\infty<s,\ b\in B\}$ (Definition~\ref{bpP:p11a:domain-def}, written $\Ddom rsB$ here), $E+t=\{b\in\C^N:\exists b'\in E,\ \lvert b-b'\rvert_\infty<t\}$, and $\nrmm F{r,s;B}$ the supremum. Euclidean norms and the Euclidean real thickening $K_\vartheta=\{y\in\R^N:\operatorname{dist}_2(y,K)\le\vartheta\}$ appear only where File~11a-1 uses them (derivative tensors $\dgm\cdot$, $\lVert\cdot\rVert^\mu_{\mathrm{dg},A}$), with the factors $\sqrt N$ stated. $N$ is the number of actions (equal to the number of angles and to the dimension of the frequency vector: $m=N$ in File~11a-1).

\begin{assumption}[Data of the selection theorem]\label{bpP:p11a:sel-data}
\begin{enumerate}[label=(\alph*)]
\item $N,\mu\ge1$ integers, $\rho\in\{1,2\}$, $\tau\ge0$ real with $\tau>\frac{N\mu}\rho-1$ (the Diophantine exponent, called $\tau_D$ in the plan).
\item $K\subset\R^N$ compact with $\operatorname{meas}K>0$; $\vartheta\in(0,1]$; $s\in(0,1]$ (angular strip width); $C_1\ge1$; $\varepsilon^*\in(0,\operatorname{meas}K)$. (Normalisation: $C_1$ occurs only as an upper bound for $\lVert D\omega\rVert$, so it may always be enlarged to $\max\{1,C_1\}$; this is done so that $\varepsilon_{\mathrm{adm}}$ is not made small by a small Hessian bound, review P11A3-01.)
\item (Hamiltonian) $U:=K+2\vartheta\subset\C^N$. $h:U\to\C$ holomorphic, real on $U\cap\R^N$; $\omega:=\nabla h$ and $\lVert D\omega(b)\rVert=\lvert D^2h(b)\rvert_{\infty\to\infty}\le C_1$ for all $b\in U$. $f:T_s\times U\to\C$ holomorphic, $1$-periodic in $\theta$, real on $\R^N\times(U\cap\R^N)$, with $\lvert f\rvert\le\varepsilon$ on $T_s\times U$.
\item (finite type, File~11a-1) $V:=U\cap\R^N\supset K_\vartheta$ is open in $\R^N$, $\omega\in C^\infty(V,\R^N)$ and the hypotheses \ref{bpP:p11a:D1}, \ref{bpP:p11a:D2} of Theorem~\ref{bpP:p11a:diophantine-measure} hold with $m=N$: $\beta_{\rho,\mu}(\omega,K)\ge\beta>0$ and $\opnm{D^kg_\rho[\omega](c,y)}\le M$ ($1\le k\le\mu+1$, $y\in K_\vartheta$), $M\ge\beta$. Put $M_0:=\lVert\omega\rVert^\mu_{\mathrm{dg},K_\vartheta}$.
\end{enumerate}
The loss $\gamma$: $\gamma:=\gamma_*(\varepsilon^*)/2$ with $\gamma_*(\varepsilon^*)$ of Corollary~\ref{bpP:p11a:power-law} (with $m=N$ and this $\tau$); thus $2\gamma\le\gamma_*(\varepsilon^*)$.
\end{assumption}
Only $\lVert D\omega\rVert\le C_1$ (with $C_1\ge1$) is used from the plan's ``$C_1\ge\lvert\omega\rvert,\lvert D^2h\rvert$''. The theorem has \emph{no} hypothesis on $\varepsilon$ beyond $\varepsilon\le\varepsilon_{\mathrm{adm}}$ (Definition~\ref{bpP:p11a:adm-params}), and no hypothesis that any conclusion holds.

\paragraph{External facts.} Besides the facts of Files~11a-1, 11a-2 and \Kref{bp:p1:ext-mvt}, \Kref{bp:p1:ext-cauchy}:
\begin{enumerate}[label=(Y\arabic*),ref=(Y\arabic*)]
\item\label{bpP:p11a:y-log} \emph{Tangent line of the logarithm.} $\log x\le\log y+x/y-1$ for $x,y>0$ (concavity; Mathlib: \nolinkurl{Real.log_le_sub_one_of_pos} applied to $x/y$).
\item\label{bpP:p11a:y-conv} \emph{Convolution with a compactly supported $C^k$ kernel.} If $\psi\in C^k_c(\R^N)$ and $g$ is bounded measurable with compact support, then $g*\psi\in C^k$ and $\partial^\alpha(g*\psi)=g*\partial^\alpha\psi$, $\lvert\partial^\alpha(g*\psi)\rvert\le\lVert g\rVert_\infty\lVert\partial^\alpha\psi\rVert_{L^1}$ (differentiation under the integral; Mathlib: \nolinkurl{HasCompactSupport.hasFDerivAt_convolution_right}).
\item\label{bpP:p11a:y-down} \emph{Continuity from above.} For a decreasing sequence of Lebesgue measurable sets of finite measure, $\operatorname{meas}\bigcap_\nu A_\nu=\lim_\nu\operatorname{meas}A_\nu$ (\nolinkurl{MeasureTheory.tendsto_measure_iInter_atTop}).
\item\label{bpP:p11a:y-geom} Geometric series $\sum_{j\ge m}x^j=x^m/(1-x)$, $0\le x<1$, and $1/(1-x)\le2$ for $x\le\frac12$.
\end{enumerate}

\paragraph{Lean totality.} Divisions are by $\gamma,T_\nu,s,s_\nu,\sigma_\nu,t_\nu,D_\nu,C_1,\varpi,q$, all positive by construction; $\log$ is applied to positive numbers ($16c_N/\varpi$, $32N/s$, $\pi T_\nu$); powers $x^{\tau}$ have $x\ge1$ or $x>0$ with real exponent as in Definition~\ref{bpP:p0:diophantine}.

%%%%%%%%%%%%%%%%%%%%%%%%%%%%%%%%%%%%%%%%%%%%%%%%%%%%%%%%%%%%%%%%%%%%%%%%%%%%%%%%
\subsection{Geometric recursion and the constants of the schedule}\label{bpP:p11a:sec-consts}

Track~K's \Kref{bp:p1:iteration} is a \emph{quadratic} recursion $\varepsilon_{\nu+1}\le c\,\varepsilon_\nu^2/(\ldots)$ and does not apply: here the one-step bound of Proposition~\ref{bpP:p11a:canonical-step}(iii) has a term linear in $\varepsilon_\nu$ (cutoff errors $2q^2+2E_T$) beside the quadratic one. The version needed is the following.

\begin{lemma}[Geometric recursion]\label{bpP:p11a:geom-recursion}
Let $\varpi\in(0,1)$, $c_0\in[0,\frac38\varpi]$, $p>0$ with $p\varpi\le1$, $P_0>0$, $x_0\ge0$ with $P_0x_0\le\frac12\varpi$, and let $x_\nu\ge0$ satisfy
\[x_{\nu+1}\ \le\ x_\nu\bigl(c_0+P_0p^{\nu}x_\nu\bigr)\qquad(\nu\ge0).\]
Then for all $\nu\ge0$: $x_\nu\le x_0\varpi^\nu$ and $P_0p^\nu x_\nu\le\frac12\varpi$. In particular $x_{\nu+1}\le\frac78\varpi x_\nu$.
\source{Ours; replaces \Kref{bp:p1:iteration} for the geometric scheme (Rus01 (13.1.7): $\varepsilon_\nu=\varepsilon_0\delta^{\kappa\nu}$, p.~156).}
\end{lemma}
\begin{proof}
\proofstep{Induction}\label{stepP:p11a:gr-ind} Suppose $x_\nu\le x_0\varpi^\nu$. Then $P_0p^\nu x_\nu\le P_0x_0(p\varpi)^\nu\le\frac12\varpi$ because $p\varpi\le1$. Hence $x_{\nu+1}\le x_\nu(\frac38\varpi+\frac12\varpi)=\frac78\varpi x_\nu\le\varpi x_\nu\le x_0\varpi^{\nu+1}$. The case $\nu=0$ is the hypothesis $P_0x_0\le\frac12\varpi$.
\end{proof}

\begin{definition}[Constants of the schedule; depend only on $(N,\mu,\tau)$]\label{bpP:p11a:schedule-consts}
\[E_*=2(\tau+1)(\mu+1)+N+4,\quad \varpi=2^{-2E_*},\quad q=\tfrac14\sqrt\varpi=2^{-E_*-2},\]
\[\alpha=\tfrac18\varpi,\qquad \varrho=\varpi/q=4\sqrt\varpi=2^{2-E_*},\]
\[\Theta=2,\qquad\zeta=\tfrac12,\qquad C_Q=18N(4N+5),\qquad A_0=\log(16c_N/\varpi)\]
($c_N$ of Lemma~\ref{bpP:p11a:fourier-cutoff}). Further
\[g_1=\varrho\,\Theta^{\tau}\zeta^{-(N+1)},\qquad g_2=\varrho\,\Theta^{\tau+1},\qquad g_3=\varrho\,\Theta^{(\tau+1)\mu},\qquad g_4=g_1\,\zeta\,\Theta^{\tau+1}.\]
\source{Ours. $\varpi$ replaces Rus01's $\delta^\kappa$ ($\kappa=24/25$, $\tau_0=1/9$ in (13.1), p.~156), $\Theta=2$ replaces the coupling $T_\nu=\Psi^{-1}(\Psi(T_0)\delta^{\tau\nu})$ of 13.1.4; deviation D-16 (Section~\ref{bpP:p11a:notes3}).}
\end{definition}

\begin{lemma}[The growth factors are beaten]\label{bpP:p11a:schedule-ratios}
With the constants of Definition~\ref{bpP:p11a:schedule-consts}: $q\le2^{-5}$, $2q^2=\frac18\varpi$, $2\alpha=\frac14\varpi$, $\varrho\le2^{-5}2^{-(\tau+1)}$ and
\[g_1,\ g_2,\ g_3,\ g_4\ \le\ 2^{-5}.\]
Consequently $q\le\zeta$, $\varrho\le\frac12$ and $\varrho\,\Theta^{\tau+1}\le\frac12$.
\source{Ours; this is where the exponent conditions of Rus01 13.1.2 ($\kappa-\sigma-\tau_0>\ldots$) and the summability $\tau\mu<\tau_0$ of (20.33) are paid.}
\end{lemma}
\begin{proof}
\proofstep{Exponents}\label{stepP:p11a:sr-exp} $\varrho=2^{2-E_*}$, so $g_1=2^{2-E_*+\tau+N+1}$ and $E_*-\tau-N-3=2(\tau+1)(\mu+1)-\tau+1\ge4(\tau+1)-\tau+1=3\tau+5\ge5$ ($\mu\ge1$). $g_2=2^{2-E_*+\tau+1}$ and $E_*-\tau-3\ge3\tau+N+5\ge5$. $g_3=2^{2-E_*+(\tau+1)\mu}$ and $E_*-2-(\tau+1)\mu=(\tau+1)\mu+2(\tau+1)+N+2\ge6$. $g_4=2^{2\tau+N+3-E_*}$ and $E_*-2\tau-N-3=2(\tau+1)(\mu+1)+1-2\tau\ge2\tau+5\ge5$. Finally $E_*\ge9$ (equality for $(N,\mu,\tau)=(1,1,0)$, which is permitted), so $q=2^{-E_*-2}\le2^{-11}\le2^{-5}$, $2q^2=\varpi/8$ and $2\alpha=\varpi/4$ by definition; $\varrho\,2^{\tau+1}=g_2\le2^{-5}$.
\end{proof}

%%%%%%%%%%%%%%%%%%%%%%%%%%%%%%%%%%%%%%%%%%%%%%%%%%%%%%%%%%%%%%%%%%%%%%%%%%%%%%%%
\subsection{The schedule}\label{bpP:p11a:sec-schedule}

\begin{definition}[The schedule]\label{bpP:p11a:schedule}
Let the constants of Definition~\ref{bpP:p11a:schedule-consts} and the data of Assumption~\ref{bpP:p11a:sel-data} be given, and let $T_0\ge1$, $t_0\in(0,1]$, $s_0>0$, $\varepsilon_0>0$ be numbers (chosen in Definition~\ref{bpP:p11a:adm-params}); put $\varsigma_0=s/24$ and $\lambda_0=\varepsilon_0/s_0$. For $\nu\ge0$:
\[T_\nu=T_02^\nu,\qquad D_\nu=\gamma T_\nu^{-\tau},\qquad t_\nu=t_02^{-(\tau+1)\nu},\qquad\varepsilon_\nu=\varepsilon_0\varpi^\nu,\qquad s_\nu=s_0q^\nu,\]
\[\delta_\nu=\sigma_\nu=\varsigma_02^{-\nu},\qquad d_\nu=d_{T_\nu}(\alpha)=\frac{\log(2c_N/\alpha)+N\log(\pi T_\nu)}{\pi T_\nu},\]
\[r_0=s,\qquad r_{\nu+1}=r_\nu-d_\nu-\delta_\nu-2\sigma_\nu,\]
\[\ell_\nu=2\varepsilon_\nu/s_\nu,\qquad S_m=\sum_{j\ge m}\ell_j,\qquad E_\nu=\varepsilon_N(T_\nu,d_\nu)\quad(\text{Lemma~\ref{bpP:p11a:fourier-cutoff}}),\qquad Q_\nu^\sharp=\frac{C_Q(2/\delta_\nu)^N\varepsilon_\nu}{\pi D_\nu s_\nu\sigma_\nu}.\]
The \emph{shells} are $\Sh_\nu=\{k\in\Z^N:\ T_{\nu-1}<\lvert k\rvert_1\le T_\nu\}$ with $T_{-1}:=0$ (so $\Sh_0=\{0<\lvert k\rvert_1\le T_0\}$), and $\nu(k)$ is the unique $\nu$ with $k\in\Sh_\nu$ ($k\ne0$); $\lvert k\rvert_1\le T_{\nu(k)}$. The $\ell^1$ norm is the one of the cutoff (File~11a-2 truncates in $\lvert k\rvert_1$).
\end{definition}

\begin{lemma}[The strip budget]\label{bpP:p11a:strip-budget}
Let $T_0\ge1$ and $T_0\ge\frac{16}{\pi s}\bigl(A_0+N\log\frac{32N}s\bigr)$ (assumption (A6) below). Then
\[\sum_{\nu\ge0}d_\nu\le\frac s4,\quad d_\nu\le\frac s4\le\frac14,\quad\sum_{\nu\ge0}(\delta_\nu+2\sigma_\nu)=\frac s4,\quad r_\nu\ge\frac s2,\quad r_{\nu+1}=r_\nu-d_\nu-\delta_\nu-2\sigma_\nu\ge\frac s2,\]
and $E_\nu\le\alpha$ for all $\nu$ (so hypothesis (C1) of Proposition~\ref{bpP:p11a:canonical-step} holds at every stage).
\source{Ours; replaces Rus01 13.1.5--6 and Lemma 13.2.2 (the Bruno-type integral condition $\tau\log\frac1\delta\int_{T_0}^\infty\log\frac1{\Phi}\frac{dT}{T^2}\le\ldots$ and $\sum1/T_\mu$), review A3 items 2, 3; deviation D-15.}
\end{lemma}
\begin{proof}
\proofstep{Sum of the cutoff losses}\label{stepP:p11a:sb-sum} Lemma~\ref{bpP:p11a:loss-sum} with $\Theta=2$, $\alpha_\nu=\alpha$ ($a=\log(2c_N/\alpha)=A_0$, $b=0$): $\sum_\nu d_\nu\le\frac1{\pi T_0}[(A_0+N\log(\pi T_0))\cdot2+N\log2\cdot2]=\frac4X(A_0+N\log X)$ with $X=2\pi T_0$.

\proofstep{Solving for $T_0$}\label{stepP:p11a:sb-solve} By \ref{bpP:p11a:y-log} with $y=32N/s$: $N\log X\le\frac{sX}{32}+N\log\frac{32N}s$, so $\frac4X(A_0+N\log X)\le\frac s8+\frac4X(A_0+N\log\frac{32N}s)$. The assumption on $T_0$ reads $X\ge\frac{32}s(A_0+N\log\frac{32N}s)$, i.e.\ $\frac4X(A_0+N\log\frac{32N}s)\le\frac s8$. Hence $\sum d_\nu\le\frac s4$; each $d_\nu\le\sum d_j\le\frac s4\le\frac14$.

\proofstep{Losses of $\delta_\nu,\sigma_\nu$ and the strip}\label{stepP:p11a:sb-strip} $\sum_\nu(\delta_\nu+2\sigma_\nu)=3\varsigma_0\sum2^{-\nu}=6\varsigma_0=\frac s4$. Therefore $r_\nu=s-\sum_{j<\nu}(d_j+\delta_j+2\sigma_j)\ge s-\frac s4-\frac s4=\frac s2$, for every $\nu$ (including $\nu+1$ in the last identity).

\proofstep{Cutoff error}\label{stepP:p11a:sb-err} $d_\nu=d_{T_\nu}(\alpha)\le1$ and $T_\nu\ge1$, $\alpha\le1$: Lemma~\ref{bpP:p11a:loss-choice} gives $E_\nu=\varepsilon_N(T_\nu,d_{T_\nu})\le\alpha$.
\end{proof}

\begin{lemma}[Elementary properties of the schedule]\label{bpP:p11a:schedule-props}
Let $\lambda_0=\varepsilon_0/s_0$, $a=\gamma T_0^{-\tau-1}$ and suppose $4\lambda_0\le a$, $8\lambda_0\le C_1t_0$ and $4C_1t_0\le a$ (these are (A1), (A2) below). For all $\nu,m\ge0$:
\begin{enumerate}[label=(S\arabic*),ref=(S\arabic*)]
\item\label{bpP:p11a:SC1} $T_{\nu+1}=2T_\nu$; $t_{\nu+1}=2^{-(\tau+1)}t_\nu\le\frac12t_\nu$; $t_\nu\le\gamma T_\nu^{-\tau-1}/(4C_1)$; $s_\nu=s_0q^\nu$ and $\varepsilon_\nu/s_\nu=\lambda_0\varrho^\nu$, $\ell_\nu=2\lambda_0\varrho^\nu$.
\item\label{bpP:p11a:SC2} (tail) $S_m\le4\lambda_0\varrho^m\le\gamma T_m^{-\tau-1}$.
\item\label{bpP:p11a:SC3} (derivatives of the increments) $\sum_{j<\nu}\dfrac{4\varepsilon_j}{s_jt_j}\le\dfrac{8\lambda_0}{t_0}\le C_1$, and $\sum_{j<\nu}\ell_jt_j^{-\mu}\le4\lambda_0t_0^{-\mu}$.
\item\label{bpP:p11a:SC4} (quadratic coefficient) $Q^\sharp_\nu=Q_0g_1^\nu$ with $Q_0:=\dfrac{C_Q2^NT_0^\tau\lambda_0}{\pi\gamma\varsigma_0^{N+1}}$.
\item\label{bpP:p11a:SC5} (divisor and collar) $t_{\nu+1}T_\nu\cdot2C_1\le D_\nu/2$.
\end{enumerate}
\source{Ours.}
\end{lemma}
\begin{proof}
\proofstep{(S1)}\label{stepP:p11a:sp-1} Definition~\ref{bpP:p11a:schedule}; $2^{-(\tau+1)}\le\frac12$ as $\tau\ge0$; $t_\nu=t_02^{-(\tau+1)\nu}\le\frac a{4C_1}2^{-(\tau+1)\nu}=\gamma T_\nu^{-\tau-1}/(4C_1)$ by $4C_1t_0\le a$. $\varepsilon_\nu/s_\nu=\frac{\varepsilon_0\varpi^\nu}{s_0q^\nu}=\lambda_0\varrho^\nu$.

\proofstep{(S2)}\label{stepP:p11a:sp-2} $\varrho\le\frac12$, so $S_m=2\lambda_0\varrho^m/(1-\varrho)\le4\lambda_0\varrho^m$ (\ref{bpP:p11a:y-geom}); $\varrho\le2^{-(\tau+1)}$ (Lemma~\ref{bpP:p11a:schedule-ratios}) and $4\lambda_0\le a$ give $S_m\le a2^{-(\tau+1)m}=\gamma T_m^{-\tau-1}$.

\proofstep{(S3)}\label{stepP:p11a:sp-3} $\frac{\varepsilon_j}{s_jt_j}=\frac{\lambda_0}{t_0}g_2^j$ with $g_2\le\frac12$, so $\sum_j\frac{4\varepsilon_j}{s_jt_j}\le\frac{8\lambda_0}{t_0}\le C_1$. And $\ell_jt_j^{-\mu}=2\lambda_0t_0^{-\mu}g_3^j$, $g_3\le\frac12$.

\proofstep{(S4)}\label{stepP:p11a:sp-4} $Q^\sharp_\nu=\frac{C_Q2^N\varsigma_0^{-N}2^{N\nu}\varepsilon_0\varpi^\nu\,T_0^\tau2^{\tau\nu}}{\pi\gamma\,s_0q^\nu\,\varsigma_02^{-\nu}}=\frac{C_Q2^NT_0^\tau\lambda_0}{\pi\gamma\varsigma_0^{N+1}}\cdot\bigl(\varrho2^{N+1+\tau}\bigr)^\nu$, and $\varrho\,2^{N+1+\tau}=g_1$ ($\Theta^\tau\zeta^{-(N+1)}=2^{\tau+N+1}$).

\proofstep{(S5)}\label{stepP:p11a:sp-5} By (S1), $t_{\nu+1}T_\nu2C_1\le2^{-(\tau+1)}\frac{\gamma T_\nu^{-\tau-1}}{4C_1}T_\nu2C_1=2^{-(\tau+2)}\gamma T_\nu^{-\tau}\le D_\nu/2$.
\end{proof}

%%%%%%%%%%%%%%%%%%%%%%%%%%%%%%%%%%%%%%%%%%%%%%%%%%%%%%%%%%%%%%%%%%%%%%%%%%%%%%%%
\subsection{Mollified cutoffs and the chain extension of the frequency}\label{bpP:p11a:sec-chain}

The measure estimate of File~11a-1 tests Fourier mode $k$ against a map $\tilde\omega_k\in C^\mu$ defined on a \emph{fixed} real neighbourhood of $K$, whereas $\omega_\nu$ lives only on the shrinking complex collar $B_\nu$ of $K_\nu$. The two are connected by extending $\omega_\nu$ from $K_\nu$ (Rus01 \S20: ``chains'', Lemma 20.3--20.5, pp.~201--203; \S\S19--20.4 not read, and not followed: deviation D-7).

\begin{lemma}[Cutoff function with $t$-scaled derivatives]\label{bpP:p11a:smooth-cutoff}
Let $\mu\ge1$, $\psi_1(x)=\kappa_\mu(1-x^2)_+^{\mu+1}$ with $\kappa_\mu>0$ such that $\int\psi_1=1$ (so $\psi_1\in C^\mu(\R)$, $\operatorname{supp}\psi_1=[-1,1]$), $m_j=\lVert\psi_1^{(j)}\rVert_{L^1}$ ($0\le j\le\mu$, $m_0=1$), and
\[c_k=N^k\max_{\alpha\in\N^N,\ \lvert\alpha\rvert=k}\ \prod_{i=1}^Nm_{\alpha_i}\quad(0\le k\le\mu),\]
\[c'_{N,\mu}=\sqrt N\max_{0\le k\le\mu}\ 4^k\sum_{j=0}^k\binom kj\,c_j\,(k-j)^{k-j}\quad(0^0=1).\]
For every nonempty compact $E\subset\R^N$ and $0<t\le1$ there is $\chi=\chi_{E,t}\in C^\mu(\R^N,[0,1])$ with $\chi=1$ on $E$, $\operatorname{supp}\chi\subset\{y:\operatorname{dist}_\infty(y,E)\le t/2\}$, and for $k\le\mu$, $\lvert a_i\rvert_\infty\le1$: $\lvert D^k\chi(y)(a_1,\dots,a_k)\rvert\le c_k(4/t)^k$.
\source{Ours; replaces the smoothing of Rus01 \S19 (not read).}
\end{lemma}
\begin{proof}
\proofstep{Construction}\label{stepP:p11a:sc-constr} $a=t/4$, $E_a=\{y:\operatorname{dist}_\infty(y,E)\le a\}$ (compact), $\psi(x)=\prod_i\psi_1(x_i)\in C^\mu_c$, $\psi_a(x)=a^{-N}\psi(x/a)$ (support in $\{\lvert x\rvert_\infty\le a\}$, $\int\psi_a=1$), $\chi=\mathbf 1_{E_a}*\psi_a$. Then $0\le\chi\le1$. If $y\in E$ then $y-x\in E_a$ for all $\lvert x\rvert_\infty\le a$, so $\chi(y)=1$. If $\operatorname{dist}_\infty(y,E)>2a$ then $y-x\notin E_a$ for $\lvert x\rvert_\infty\le a$, so $\chi(y)=0$.

\proofstep{Derivatives}\label{stepP:p11a:sc-deriv} By \ref{bpP:p11a:y-conv}, $\chi\in C^\mu$ and $\lvert\partial^\alpha\chi\rvert\le\lVert\partial^\alpha\psi_a\rVert_{L^1}=a^{-\lvert\alpha\rvert}\prod_im_{\alpha_i}$ (scaling and Fubini). For $\lvert a_i\rvert_\infty\le1$, $\lvert D^k\chi(y)(a_1,\dots,a_k)\rvert\le\sum_{i_1,\dots,i_k}\lvert\partial_{i_1}\cdots\partial_{i_k}\chi\rvert\le N^k\max_{\lvert\alpha\rvert=k}\lvert\partial^\alpha\chi\rvert\le c_ka^{-k}=c_k(4/t)^k$.
\end{proof}

\begin{lemma}[Chain extension and closeness]\label{bpP:p11a:finite-extension}
Let Assumption~\ref{bpP:p11a:sel-data} and a schedule (Definition~\ref{bpP:p11a:schedule}) with $t_0\le\vartheta/\sqrt N$ be given, $\nu\ge0$, $K=K_0\supset K_1\supset\dots\supset K_\nu$ nonempty compact, $B_j=K_j+t_j$, and for $j<\nu$ let $\Delta\omega_j:B_j\to\C^N$ be holomorphic, real on $B_j\cap\R^N$, with $\lvert\Delta\omega_j\rvert_\infty\le\ell_j$. Put $\omega_\nu=\omega+\sum_{j<\nu}\Delta\omega_j$ on $B_\nu$ and, with $\chi_j=\chi_{K_j,t_j}$ of Lemma~\ref{bpP:p11a:smooth-cutoff},
\[\tilde\omega_\nu\ :=\ \omega+\sum_{j<\nu}\chi_j\,\Delta\omega_j\quad\text{on }V\qquad(\chi_j\Delta\omega_j:=0\text{ outside }\operatorname{supp}\chi_j).\]
Then:
\begin{enumerate}[label=(\alph*)]
\item $\tilde\omega_\nu\in C^\mu(V,\R^N)$ and $\tilde\omega_\nu=\omega_\nu$ on $K_\nu\subset K_j$ (all $j\le\nu$); $\tilde\omega_0=\omega$.
\item $\lVert\tilde\omega_\nu-\omega\rVert^\mu_{\mathrm{dg},K_\vartheta}\le c'_{N,\mu}\sum_{j<\nu}\ell_jt_j^{-\mu}$.
\end{enumerate}
\source{Ours (deviation D-7); the role of Rus01 (20.7), (20.25), (20.26), pp.~201--203, where the closeness $\lVert\tilde\chi_\nu-\chi\rVert^{\mu_0}\le\frac12\alpha^{4p^2}\beta/M^*$ is the hypothesis of Thm 18.5.}
\end{lemma}
\begin{proof}
\proofstep{Supports}\label{stepP:p11a:fe-supp} $\operatorname{supp}\chi_j\subset\{\operatorname{dist}_\infty(\cdot,K_j)\le t_j/2\}\subset K+\frac{\vartheta}{2\sqrt N}$ (closed), whose Euclidean distance to $K$ is $\le\vartheta/2$: inside $K_\vartheta\subset V$ and inside $B_j\cap\R^N$ with room ($K_j+t_j/2\subset B_j$).

\proofstep{Regularity and equality}\label{stepP:p11a:fe-reg} On $B_j\cap\R^N$, $\Delta\omega_j$ is real-analytic; $\chi_j\in C^\mu$ is supported where $\Delta\omega_j$ is defined, so $\chi_j\Delta\omega_j\in C^\mu(V,\R^N)$. On $K_\nu\subset K_j$ ($j<\nu$) $\chi_j=1$, so $\tilde\omega_\nu=\omega+\sum_{j<\nu}\Delta\omega_j=\omega_\nu$.

\proofstep{Derivatives of the increments}\label{stepP:p11a:fe-cauchy} For $y\in\R^N$ with $\operatorname{dist}_\infty(y,K_j)\le t_j/2$ we have $y\in K_j+\frac34t_j$. Lemma~\ref{bpP:p11a:param-cauchy} with $E=K_j$, $t=t_j$, $t'=\frac34t_j$ gives $\lVert D^u\Delta\omega_j(y)\rVert\le u^u\ell_j(t_j/4)^{-u}=u^u4^u\ell_jt_j^{-u}$ for $u\ge1$ (max norms on the arguments).

\proofstep{Leibniz}\label{stepP:p11a:fe-leib} For a unit vector $a\in\R^N$ ($\lvert a\rvert_2=1$, so $\lvert a\rvert_\infty\le1$) and $k\le\mu$: $D^k(\chi_j\Delta\omega_j)(a^k)=\sum_{u=0}^k\binom kuD^{k-u}\chi_j(a^{k-u})\,D^u\Delta\omega_j(a^u)$, hence in the max norm $\le\sum_u\binom ku c_{k-u}(4/t_j)^{k-u}u^u4^u\ell_jt_j^{-u}$ (for $u=0$ the bound is $\ell_j$, $0^0=1$), and the Euclidean norm of the value is at most $\sqrt N$ times the max norm. Reindexing $u\leftrightarrow k-u$ gives $\le c'_{N,\mu}\ell_jt_j^{-k}\le c'_{N,\mu}\ell_jt_j^{-\mu}$ (as $t_j\le1$). Sum over $j<\nu$.
\end{proof}

\begin{remark}[Verdict on D-7]\label{bpP:p11a:d7-verdict}
The claim of the plan is true, with the following precise form. The increments $\Delta\omega_j$ are \emph{holomorphic on the whole collar} $B_j$, so a $C^\mu$ cutoff at scale $t_j$ suffices and Whitney-type extension is not needed. The $C^\mu$-size of the $j$-th increment is $\lesssim c'\ell_jt_j^{-\mu}=2c'\lambda_0t_0^{-\mu}g_3^j$ and is summable exactly because $g_3=\varrho\,2^{(\tau+1)\mu}<1$ (it is $\le2^{-5}$ by Lemma~\ref{bpP:p11a:schedule-ratios}); this is Rus01's $\tau\mu<\tau_0$ (Rus01 16.3--16.4, (20.33)). The cost is the factor $t_0^{-\mu}$ in the closeness inequality (A5), where $t_0\lesssim\gamma T_0^{-\tau-1}/C_1$: this is the origin of the large degree in the budget (Section~\ref{bpP:p11a:sec-num}).
\end{remark}

%%%%%%%%%%%%%%%%%%%%%%%%%%%%%%%%%%%%%%%%%%%%%%%%%%%%%%%%%%%%%%%%%%%%%%%%%%%%%%%%
\subsection{The initial remainder and the admissible parameters}\label{bpP:p11a:sec-init}

\begin{lemma}[Initialisation; Rus01 \S14]\label{bpP:p11a:initial}
Under Assumption~\ref{bpP:p11a:sel-data}(c) let $0<t_0\le\vartheta$, $0<s_0\le\vartheta$, $B_0=K+t_0$ and
\[H_0(\theta,\eta,b)=h(b+\eta)+f(\theta,b+\eta)-h(b),\qquad N_0(\eta,b)=\omega(b)\cdot\eta,\qquad R_0=H_0-N_0 .\]
Then $R_0$ is holomorphic and $1$-periodic on $\Ddom s{s_0}{B_0}$, real at real points, and
\[\nrmm{R_0}{s,s_0;B_0}\ \le\ \varepsilon+\tfrac N2C_1s_0^2\ =:\ \varepsilon_0 .\]
(The translation $(\theta,\eta)\mapsto(\theta,b+\eta)$ is symplectic for each $b$; the constant $h(b)$ is removed; $h_0:=0$.)
\source{Rus01 \S14.1, (14.2)--(14.5) pp.~165, and Lemma 14.5 p.~168 (there $s_0^2\approx\vartheta\varepsilon_0/C_\bullet$, (14.13)); constants ours, in the max norm; the factor $N$ is $\lvert\eta^{\top}M\eta\rvert\le N\lVert M\rVert\lvert\eta\rvert_\infty^2$.}
\end{lemma}
\begin{proof}
\proofstep{Domain}\label{stepP:p11a:in-dom} For $b\in B_0$ pick $b'\in K$ with $\lvert b-b'\rvert<t_0\le\vartheta$; for $\lvert\eta\rvert_\infty<s_0\le\vartheta$ and $\lambda\in[0,1]$, $b+\lambda\eta\in K+2\vartheta=U$. So $h(b+\eta)$, $f(\theta,b+\eta)$ are defined and holomorphic; periodicity and reality are inherited.

\proofstep{Taylor}\label{stepP:p11a:in-taylor} $h(b+\eta)-h(b)-\omega(b)\cdot\eta=\int_0^1(1-\lambda)\,\eta^{\top}D\omega(b+\lambda\eta)\eta\,d\lambda$ and $\lvert\eta^{\top}M\eta\rvert\le\lvert\eta\rvert_\infty\lvert M\eta\rvert_1\le N\lvert\eta\rvert_\infty\lvert M\eta\rvert_\infty\le N\lVert M\rVert\lvert\eta\rvert_\infty^2$. So the quadratic part is $\le\frac N2C_1s_0^2$ and $\lvert f\rvert\le\varepsilon$. Thus $R_0=[h(b+\eta)-h(b)-\omega\cdot\eta]+f(\theta,b+\eta)$.
\end{proof}

\begin{definition}[Admissible parameters]\label{bpP:p11a:adm-params}
Under Assumption~\ref{bpP:p11a:sel-data}, with $\varsigma_0=s/24$, $a=\gamma T_0^{-\tau-1}$ and the constants of Definitions~\ref{bpP:p11a:schedule-consts} and Lemma~\ref{bpP:p11a:smooth-cutoff}, put
\[A_0=\log\frac{16c_N}\varpi,\qquad T_0=\max\Bigl\{1,\ \frac{16}{\pi s}\Bigl(A_0+N\log\frac{32N}s\Bigr)\Bigr\},\qquad t_0=\min\Bigl\{\frac\vartheta{\sqrt N},\ \frac a{4C_1}\Bigr\},\]
\[\Xi=2^{\tau+1}\Bigl(1+\frac{t_0}{\varsigma_0}+\frac{\varsigma_0}{t_0}+T_0\varsigma_0\Bigr),\qquad\theta_{\max}=\begin{cases}\beta/2&\rho=1,\\ \min\bigl\{\beta,\ \beta/(2^{\mu+1}(2M_0+\beta))\bigr\}&\rho=2,\end{cases}\]
\[\lambda_{\max}=\min\Bigl\{\frac{C_1t_0}8,\ \ \min\Bigl\{\frac\varpi2,\frac{2^{-10}}\Xi\Bigr\}\frac{\pi\gamma\varsigma_0^{N+1}}{C_Q2^NT_0^\tau},\ \ \frac{\theta_{\max}t_0^\mu}{4c'_{N,\mu}}\Bigr\},\qquad s_0=\frac{\lambda_{\max}}{NC_1},\qquad\varepsilon_{\mathrm{adm}}=\frac{\lambda_{\max}^2}{2NC_1}.\]
For $\varepsilon\le\varepsilon_{\mathrm{adm}}$ let $\varepsilon_0=\varepsilon+\frac N2C_1s_0^2$ (Lemma~\ref{bpP:p11a:initial}) and $\lambda_0=\varepsilon_0/s_0$. All quantities depend on the data $(N,\mu,\rho,\tau,\vartheta,s,C_1,\beta,M_0,\varepsilon^*)$ and on $d=\operatorname{diam}_\infty K$ and $M$ through $\gamma$ only.
\source{Ours (deviation D-9: explicit $\varepsilon_{\mathrm{adm}}$, replacing Rus01 (13.4)--(13.7), (14.19)--(14.22), pp.~158, 168--170, where $\varepsilon_0/s_0$ is fixed by a list of inequalities with undetermined constants $c_1,c_2,c_3,C^*$).}
\end{definition}

\begin{proposition}[The initial inequalities]\label{bpP:p11a:adm-choice}
For $\varepsilon\le\varepsilon_{\mathrm{adm}}$ the schedule of Definition~\ref{bpP:p11a:schedule} (with $T_0,t_0,s_0,\varepsilon_0$ of Definition~\ref{bpP:p11a:adm-params}) satisfies, with $a=\gamma T_0^{-\tau-1}$, $\lambda_0=\varepsilon_0/s_0$ and $Q_0=C_Q2^NT_0^\tau\lambda_0/(\pi\gamma\varsigma_0^{N+1})$:
\begin{enumerate}[label=(A\arabic*),ref=(A\arabic*)]
\item\label{bpP:p11a:A1} $0<t_0\le\vartheta/\sqrt N\le1$ and $4C_1t_0\le a$;
\item\label{bpP:p11a:A2} $8\lambda_0\le C_1t_0$ (hence $4\lambda_0\le a/8\le a$);
\item\label{bpP:p11a:A3} $s_0\le t_0/(8N)\ (\le\vartheta)$;
\item\label{bpP:p11a:A4} $Q_0\le\varpi/2$ and $Q_0\,\Xi\le2^{-10}$;
\item\label{bpP:p11a:A5} $4c'_{N,\mu}\lambda_0t_0^{-\mu}\le\theta_{\max}$;
\item\label{bpP:p11a:A6} $T_0\ge1$ and $T_0\ge\frac{16}{\pi s}(A_0+N\log\frac{32N}s)$ (Lemma~\ref{bpP:p11a:strip-budget});
\item\label{bpP:p11a:A7} $\varepsilon_0=\varepsilon+\frac N2C_1s_0^2$, $\lvert R_0\rvert\le\varepsilon_0$ on $\Ddom s{s_0}{B_0}$ (Lemma~\ref{bpP:p11a:initial}).
\end{enumerate}
In particular $\varepsilon_{\mathrm{adm}}>0$ and $s_0\le\vartheta$. (Monotonicity in the data is not claimed here; the polynomial lower bound is Proposition~\ref{bpP:p11a:adm-budget} in Section~\ref{bpP:p11a:sec-num}, and the monotone minorant is File~11a-5's business.)
\end{proposition}
\begin{proof}
\proofstep{$\lambda_0\le\lambda_{\max}$}\label{stepP:p11a:ac-lam} $\lambda_0=\varepsilon/s_0+\frac N2C_1s_0=\varepsilon NC_1/\lambda_{\max}+\frac12\lambda_{\max}\le\lambda_{\max}$ because $\varepsilon\le\lambda_{\max}^2/(2NC_1)$.

\proofstep{(A1)--(A5)}\label{stepP:p11a:ac-list} (A1) is the definition of $t_0$ ($\vartheta\le1$). Since $\lambda_0\le\lambda_{\max}$: (A2) from the first term of the minimum; $C_1t_0\le a/4$ gives $4\lambda_0\le a/8$. (A3): $s_0=\lambda_{\max}/(NC_1)\le t_0/(8N)$. (A4): $Q_0=\frac{C_Q2^NT_0^\tau\lambda_0}{\pi\gamma\varsigma_0^{N+1}}\le\min\{\varpi/2,2^{-10}/\Xi\}$ by the second term. (A5): third term. (A6): definition of $T_0$. (A7): Lemma~\ref{bpP:p11a:initial} (its hypotheses $t_0\le\vartheta$, $s_0\le\vartheta$ are (A1), (A3)).
\end{proof}

\begin{remark}[The three branches of $\varepsilon_{\mathrm{adm}}$]\label{bpP:p11a:adm-branches}
Write $\lambda_{\max}=\min\{\lambda^{(1)},\lambda^{(2)},\lambda^{(3)}\}$ for the three terms of Definition~\ref{bpP:p11a:adm-params}. Then $\varepsilon_{\mathrm{adm}}=\min\{\varepsilon^{(1)},\varepsilon^{(2)},\varepsilon^{(3)}\}$ with exactly
\[\varepsilon^{(1)}=\frac{C_1t_0^2}{128N},\qquad\varepsilon^{(2)}=\frac{\pi^2\gamma^2\varsigma_0^{2N+2}}{2NC_1C_Q^2\,4^{N}\,T_0^{2\tau}}\ \min\Bigl\{\frac\varpi2,\frac{2^{-10}}\Xi\Bigr\}^2,\qquad\varepsilon^{(3)}=\frac{\theta_{\max}^2\,t_0^{2\mu}}{32\,N\,c_{N,\mu}'^{\,2}\,C_1}.\]
All three can bind. For instance, if $t_0=\gamma T_0^{-\tau-1}/(4C_1)$ and $\Xi\simeq2^{\tau+1}\varsigma_0/t_0$ (the case $\varsigma_0/t_0$ dominates $\Xi$), then $\varepsilon^{(2)}\asymp\gamma^{4}\varsigma_0^{2N}C_1^{-3}T_0^{-4\tau-2}$, of higher degree in $\gamma$ than $\varepsilon^{(3)}$ when $\mu=1$; $\varepsilon^{(3)}$ is a multiple of $\theta_{\max}^2t_0^{2\mu}$, and for $\rho=2$, $\theta_{\max}$ is $\beta/(2^{\mu+1}(2M_0+\beta))$ or $\beta$, not universally $\beta^2/M_0^2$ after squaring. No statement about the size of $\varepsilon_{\mathrm{adm}}$ below is made from one branch alone.
\end{remark}

%%%%%%%%%%%%%%%%%%%%%%%%%%%%%%%%%%%%%%%%%%%%%%%%%%%%%%%%%%%%%%%%%%%%%%%%%%%%%%%%
\subsection{The finite-stage divisor invariant and the retained sets}\label{bpP:p11a:sec-retained}

\emph{Stage-$\nu$ data.} A \emph{chain of length $\nu$} consists of nonempty compact sets $K=K_0\supset K_1\supset\dots\supset K_\nu$ and holomorphic maps $\Delta\omega_j:B_j=K_j+t_j\to\C^N$ ($j<\nu$), real on $B_j\cap\R^N$, with $\lvert\Delta\omega_j\rvert_\infty\le\ell_j$ on $B_j$, such that for $j<\nu$
\begin{equation}\label{bpP:p11a:retained-def}
K_{j+1}=\bigl\{b\in K_j:\ \lvert k\cdot\omega_j(b)\rvert\ge2\gamma\lvert k\rvert_1^{-\tau}\ \text{ for all }k\in\Sh_j\bigr\},\qquad\omega_j=\omega+\sum_{i<j}\Delta\omega_i\ \text{ on }B_j.
\end{equation}
($K_{\nu+1}$ is defined by \eqref{bpP:p11a:retained-def} with $j=\nu$ once $\omega_\nu$ is known; it is compact, as $\omega_\nu$ is continuous on $K_\nu$.)

\begin{lemma}[Finite-stage divisor invariant]\label{bpP:p11a:finite-divisor}
Let a chain of length $\nu$ be given, let $K_{\nu+1}$ be defined by \eqref{bpP:p11a:retained-def}, and assume $S_m\le\gamma T_m^{-\tau-1}$ for all $m\le\nu$ (this is \ref{bpP:p11a:SC2}). Then for every $b\in K_{\nu+1}$ and every $k\in\Z^N$ with $0<\lvert k\rvert_1\le T_\nu$:
\begin{gather*}\lvert\omega_\nu(b)-\omega_{\nu(k)}(b)\rvert_\infty\le\gamma T_{\nu(k)}^{-\tau-1}\le\gamma\lvert k\rvert_1^{-\tau-1},\qquad\lvert k\cdot\omega_{\nu(k)}(b)\rvert\ge2\gamma\lvert k\rvert_1^{-\tau},\\ \boxed{\ \lvert k\cdot\omega_\nu(b)\rvert\ \ge\ \gamma\lvert k\rvert_1^{-\tau}\ \ge\ \gamma T_\nu^{-\tau}=D_\nu.\ }\end{gather*}
Thus the hypotheses of Lemma~\ref{bpP:p11a:divisor-wide} hold with $u_k=\omega_{\nu(k)}(b)$, $v=\omega_\nu(b)$, $T=T_\nu$, and the divisor hypothesis of (C2) in Proposition~\ref{bpP:p11a:canonical-step} holds on $E=K_{\nu+1}$ with $D=D_\nu$.
\source{Ours; this \emph{replaces} the limit-based tail requirement (I3) of Remark~\ref{bpP:p11a:consumers}, which bounded $\lvert\omega_*-\omega_{\nu(k)}\rvert$ on $K_\infty$ only and used the endpoint $T_{\nu(k)-1}$ (review P11A2-01); see Remark~\ref{bpP:p11a:i3-corrected}.}
\end{lemma}
\begin{proof}
\proofstep{Telescoping}\label{stepP:p11a:fd-tel} Let $m=\nu(k)\le\nu$ (as $\lvert k\rvert_1\le T_\nu$) and $b\in K_{\nu+1}\subset K_j\subset B_j$ for all $j\le\nu$. On $B_\nu$, $\omega_\nu-\omega_m=\sum_{j=m}^{\nu-1}\Delta\omega_j$ (both defined at $b$, since $B_\nu\subset B_m$), so $\lvert\omega_\nu(b)-\omega_m(b)\rvert_\infty\le\sum_{j=m}^{\nu-1}\ell_j\le S_m\le\gamma T_m^{-\tau-1}$; and $T_m^{-\tau-1}\le\lvert k\rvert_1^{-\tau-1}$ because $\lvert k\rvert_1\le T_m$, $\tau+1>0$.

\proofstep{Mode $k$ was tested at stage $m$}\label{stepP:p11a:fd-test} $b\in K_{\nu+1}\subset K_{m+1}$ and $k\in\Sh_m$, so by \eqref{bpP:p11a:retained-def} $\lvert k\cdot\omega_m(b)\rvert\ge2\gamma\lvert k\rvert_1^{-\tau}$.

\proofstep{Conclusion}\label{stepP:p11a:fd-concl} By $\lvert k\cdot x\rvert\le\lvert k\rvert_1\lvert x\rvert_\infty$: $\lvert k\cdot(\omega_\nu-\omega_m)(b)\rvert\le\lvert k\rvert_1\gamma\lvert k\rvert_1^{-\tau-1}=\gamma\lvert k\rvert_1^{-\tau}$, so $\lvert k\cdot\omega_\nu(b)\rvert\ge2\gamma\lvert k\rvert_1^{-\tau}-\gamma\lvert k\rvert_1^{-\tau}$. Finally $\lvert k\rvert_1\le T_\nu$, $\tau\ge0$. (This is also the conclusion of Lemma~\ref{bpP:p11a:divisor-wide}, now with all its hypotheses verified; for $m=\nu$ the telescoping sum is empty.)
\end{proof}

\begin{remark}[The correct form of (I3)]\label{bpP:p11a:i3-corrected}
Requirement \ref{bpP:p11a:I3} of Remark~\ref{bpP:p11a:consumers} should read: \emph{for every $m\ge0$ (including $m=0$)} $\sum_{j\ge m}\ell_j\le\gamma\,T_m^{-\tau-1}$, where $\ell_j\ge\sup_{B_j}\lvert\Delta\omega_j\rvert_\infty$; and then, for $k$ with $T_{m-1}<\lvert k\rvert_1\le T_m$ (i.e.\ $m=\nu(k)$), $\lvert\omega_*(b)-\omega_m(b)\rvert_\infty\le\sum_{j\ge m}\ell_j\le\gamma T_m^{-\tau-1}\le\gamma\lvert k\rvert_1^{-\tau-1}$. The printed endpoint $T_{m-1}$ is on the wrong side: it would demand the bound $\gamma\lvert k\rvert_1^{-\tau-1}$ only for the smallest $\lvert k\rvert_1$ of the shell, where it is weakest, and gives no control of the largest. In addition the finite-stage version (Lemma~\ref{bpP:p11a:finite-divisor}) must be proved for $\omega_\nu-\omega_m$ on the finite retained set, \emph{not} derived from the limit (two bounds to the limit would give $2\gamma$ instead of $\gamma$). The schedule satisfies this by \ref{bpP:p11a:SC2}.
\end{remark}

\begin{lemma}[The retained set is nonempty: finite-stage measure bound]\label{bpP:p11a:finite-measure}
Let Assumption~\ref{bpP:p11a:sel-data}, the schedule with (A1)--(A6), a chain of length $\nu$ and $K_{\nu+1}$ as in \eqref{bpP:p11a:retained-def} be given. Put
\[H^{(\nu)}=\Bigl\{b\in K:\ \lvert k\cdot\tilde\omega_{\nu(k)}(b)\rvert\ge2\gamma\lvert k\rvert_1^{-\tau}\ \text{ for all }k\in\Z^N,\ 0<\lvert k\rvert_1\le T_\nu\Bigr\}\]
with $\tilde\omega_j$ the chain extension of Lemma~\ref{bpP:p11a:finite-extension} (it uses $K_i,\Delta\omega_i$ for $i<j\le\nu$ only). Then $H^{(\nu)}$ is compact, $K_{\nu+1}\supset H^{(\nu)}$, $H^{(\nu)}\supset H^{(\nu+1)}$ whenever both are defined, and
\[\operatorname{meas}H^{(\nu)}\ \ge\ \operatorname{meas}K-\varepsilon^*\ >\ 0 .\]
In particular $K_{\nu+1}\ne\emptyset$.
\source{Ours; Rus01 Thm 20.5.1 (any maximal chain is infinite), p.~201--202, whose proof is this argument with Thm 18.5 (here Corollary~\ref{bpP:p11a:power-law}); the circularity (the sets $K_j$ depend on the construction) is resolved as there: only finitely many stages, already built, enter $H^{(\nu)}$.}
\end{lemma}
\begin{proof}
\proofstep{Closeness for every stage}\label{stepP:p11a:fm-close} For $j\le\nu$ the chain of length $j$ is an initial segment, so Lemma~\ref{bpP:p11a:finite-extension}(b) and \ref{bpP:p11a:SC3}, (A5) give $\theta_j:=\lVert\tilde\omega_j-\omega\rVert^\mu_{\mathrm{dg},K_\vartheta}\le c'_{N,\mu}\cdot4\lambda_0t_0^{-\mu}\le\theta_{\max}$. By Lemma~\ref{bpP:p11a:stability}, $\eta_\rho(\theta_j)\le\beta/2$ (for $\rho=2$: $\theta_j\le\beta$ and $2^\mu\theta_j(2M_0+\theta_j)\le2^\mu\theta_j(2M_0+\beta)\le\beta/2$ by definition of $\theta_{\max}$), so \ref{bpP:p11a:D3} holds for $\tilde\omega_j$.

\proofstep{A family indexed by all modes}\label{stepP:p11a:fm-family} For $k\ne0$ put $\tilde\omega_k:=\tilde\omega_{\min\{\nu(k),\nu\}}$. Then \ref{bpP:p11a:D1}--\ref{bpP:p11a:D3} hold for the family $(\tilde\omega_k)$. Corollary~\ref{bpP:p11a:power-law} with $2\gamma\le\gamma_*(\varepsilon^*)$ in place of $\gamma$ gives that $H_{2\gamma}:=\{b\in K:\lvert k\cdot\tilde\omega_k(b)\rvert\ge2\gamma\lvert k\rvert_1^{-\tau}\ \forall k\ne0\}$ is compact with $\operatorname{meas}H_{2\gamma}\ge\operatorname{meas}K-\varepsilon^*$. For $\lvert k\rvert_1\le T_\nu$, $\tilde\omega_k=\tilde\omega_{\nu(k)}$, so $H_{2\gamma}\subset H^{(\nu)}$.

\proofstep{Inclusion in the retained set}\label{stepP:p11a:fm-incl} Let $b\in H^{(\nu)}$. By induction $b\in K_j$ for $j\le\nu+1$: $b\in K_0$; if $b\in K_j$, $j\le\nu$, then $\omega_j(b)=\tilde\omega_j(b)$ (Lemma~\ref{bpP:p11a:finite-extension}(a)), and each $k\in\Sh_j$ has $\nu(k)=j$, $\lvert k\rvert_1\le T_j\le T_\nu$, so $\lvert k\cdot\omega_j(b)\rvert=\lvert k\cdot\tilde\omega_{\nu(k)}(b)\rvert\ge2\gamma\lvert k\rvert_1^{-\tau}$, whence $b\in K_{j+1}$ by \eqref{bpP:p11a:retained-def}. Compactness of $H^{(\nu)}$: an intersection of finitely many closed sets of $K$ (the $\tilde\omega_{\nu(k)}$ are continuous).
\end{proof}

%%%%%%%%%%%%%%%%%%%%%%%%%%%%%%%%%%%%%%%%%%%%%%%%%%%%%%%%%%%%%%%%%%%%%%%%%%%%%%%%
\subsection{The inductive lemma}\label{bpP:p11a:sec-induct}

\begin{proposition}[Inductive lemma, maximal case; Rus01 Lemma 13.4 p.~158]\label{bpP:p11a:inductive-lemma}
Let Assumption~\ref{bpP:p11a:sel-data} hold, $\varepsilon\le\varepsilon_{\mathrm{adm}}$, and take the schedule of Definition~\ref{bpP:p11a:schedule} with the parameters of Definition~\ref{bpP:p11a:adm-params}. Say that \emph{$(\mathrm P_\nu)$ holds} if there are a chain of length $\nu$ (so $K_\nu$, $B_\nu=K_\nu+t_\nu$, $\omega_\nu=\omega+\sum_{j<\nu}\Delta\omega_j$ and $\Delta h_j$, $j<\nu$) and functions $h_\nu=\sum_{j<\nu}\Delta h_j$ (with $h_0=0$), $R_\nu$ such that
\begin{enumerate}[label=(\alph*)]
\item $h_\nu,\omega_\nu$ are holomorphic on $B_\nu$ and real on $B_\nu\cap\R^N$; $\lvert\Delta h_j\rvert\le2\varepsilon_j$ on $B_j$ and $\lvert\Delta\omega_j\rvert_\infty\le\ell_j$ on $B_j$ ($j<\nu$); $\lVert D\omega_\nu\rVert\le2C_1$ on $B_\nu$;
\item $R_\nu$ is holomorphic and $1$-periodic on $\Ddom{r_\nu}{s_\nu}{B_\nu}$, real at real points, $\nrmm{R_\nu}{r_\nu,s_\nu;B_\nu}\le\varepsilon_\nu$;
\item $H_\nu:=N_\nu+R_\nu$ with $N_\nu(\eta,b)=h_\nu(b)+\omega_\nu(b)\cdot\eta$ satisfies $H_\nu=H_0\circ X_0\circ\dots\circ X_{\nu-1}$ on $\Ddom{r_\nu}{s_\nu}{B_\nu}$, where $H_0$ is that of Lemma~\ref{bpP:p11a:initial} and $X_j$ are the maps of the next statement.
\end{enumerate}
Then $(\mathrm P_0)$ holds, and for every $\nu\ge0$, $(\mathrm P_\nu)$ implies the existence of $K_{\nu+1}\ne\emptyset$ compact (by \eqref{bpP:p11a:retained-def}), $\Delta h_\nu,\Delta\omega_\nu$ and $X_\nu$ such that $(\mathrm P_{\nu+1})$ holds with $B_{\nu+1}=K_{\nu+1}+t_{\nu+1}$ and:
\begin{enumerate}[label=(O\arabic*),ref=(O\arabic*)]
\item\label{bpP:p11a:O1} (finite-stage divisors) for $b\in K_{\nu+1}$, $0<\lvert k\rvert_1\le T_\nu$: $\lvert k\cdot\omega_\nu(b)\rvert\ge\gamma\lvert k\rvert_1^{-\tau}\ge D_\nu$ (Lemma~\ref{bpP:p11a:finite-divisor});
\item\label{bpP:p11a:O2} (normal form) $\Delta h_\nu,\Delta\omega_\nu$ are holomorphic and real on $B_\nu$ with $\lvert\Delta h_\nu\rvert\le2\varepsilon_\nu$, $\lvert\Delta\omega_\nu\rvert_\infty\le\ell_\nu=2\varepsilon_\nu/s_\nu$; $\omega_{\nu+1}=\omega_\nu+\Delta\omega_\nu$, $h_{\nu+1}=h_\nu+\Delta h_\nu$ on $B_{\nu+1}$;
\item\label{bpP:p11a:O3} (the map) for each $b\in B_{\nu+1}$, $X_{\nu,b}=\Phi_b$ of Proposition~\ref{bpP:p11a:canonical-step} is a simple canonical map on $\Ddom{r_{\nu+1}}{s_{\nu+1}}{\{b\}}$, holomorphic jointly in $(\theta,\eta,b)$ on $\Ddom{r_{\nu+1}}{s_{\nu+1}}{B_{\nu+1}}$, with values in $\Ddom{r_{\nu+1}+\sigma_\nu}{s_\nu}{B_{\nu+1}}\subset\Ddom{r_\nu}{s_\nu}{B_{\nu+1}}$, $X_\nu(\theta+k,\eta)=X_\nu(\theta,\eta)+(k,0)$, real at real points; with $A_{1,\nu}=\LamD{\delta_\nu}(1+E_\nu)\varepsilon_\nu/(\pi D_\nu s_\nu)$, $A_{\eta,\nu}=(N+1)s_\nu A_{1,\nu}/\sigma_\nu$, $\kappa_\nu=2(N+1)A_{1,\nu}/\sigma_\nu$:
$\lvert X_{\nu,\theta}-\theta\rvert_\infty\le A_{1,\nu}$, $\lvert X_{\nu,\eta}-\eta\rvert_\infty\le A_{\eta,\nu}$, $\kappa_\nu\le\frac12$, and the unscaled derivative bound $\lvert DX_\nu-\mathbf 1\rvert_\infty\le2\kappa_\nu+\kappa_\nu s_\nu/\sigma_\nu$ in $(\theta,\eta)$; on the real part $\lvert D\Theta_{\nu,b}-\mathbf 1\rvert_\infty\le e^{\Lambda_{\mathbb R,\nu}}-1$ with $\Lambda_{\mathbb R,\nu}=2T_\nu(1+E_\nu)\varepsilon_\nu\LamD{r_\nu-d_\nu}/(D_\nu s_\nu)$;
\item\label{bpP:p11a:O4} (parameter derivatives) for $0<t'<t_{\nu+1}$ and all $(\theta,\eta)$ in the domain: $\lVert D_b^j(X_\nu-\mathrm{id})_\theta\rVert\le j^jA_{1,\nu}(t_{\nu+1}-t')^{-j}$ and likewise for the $\eta$-component with $A_{\eta,\nu}$ on $K_{\nu+1}+t'$;
\item\label{bpP:p11a:O5} (conjugation) $H_\nu\circ X_\nu=N_{\nu+1}+R_{\nu+1}$ on $\Ddom{r_{\nu+1}}{s_{\nu+1}}{B_{\nu+1}}$.
\end{enumerate}
\source{Rus01 Lemma 13.4, p.~158 (conditions (13.4)--(13.7) become (A1)--(A7); $\lvert\Delta\omega_\nu\rvert\le\frac32\varepsilon_\nu/s_\nu$ becomes $\le2\varepsilon_\nu/s_\nu$; the Lipschitz factor $e^{24C_\nu T_\nu}$ of 13.4.5.4 becomes $e^{\Lambda_{\mathbb R,\nu}}$, Lemma~\ref{bpP:p11a:real-lipschitz}); deviations D-7, D-9, D-15, D-16. Proof by Proposition~\ref{bpP:p11a:canonical-step}.}
\end{proposition}
\begin{proof}
\proofstep{$(\mathrm P_0)$}\label{stepP:p11a:il-p0} $K_0=K$, $B_0=K+t_0$, $\Delta\omega_j$ empty, $\omega_0=\omega$, $h_0=0$, $R_0$ of Lemma~\ref{bpP:p11a:initial}: $\lvert R_0\rvert\le\varepsilon_0$ on $\Ddom s{s_0}{B_0}=\Ddom{r_0}{s_0}{B_0}$ (\ref{bpP:p11a:A7}, $r_0=s$). $\lVert D\omega\rVert\le C_1\le2C_1$ on $B_0\subset U$ since $t_0\le\vartheta$. $H_0$ is the composition of the empty list.

\proofstep{Retained set, nonempty}\label{stepP:p11a:il-ret} Given $(\mathrm P_\nu)$ with its chain, \ref{bpP:p11a:SC2} holds by Lemma~\ref{bpP:p11a:schedule-props} (hypotheses \ref{bpP:p11a:A1}, \ref{bpP:p11a:A2}). $K_{\nu+1}$ of \eqref{bpP:p11a:retained-def} is compact and by Lemma~\ref{bpP:p11a:finite-measure} contains $H^{(\nu)}$, which has positive measure; so $K_{\nu+1}\neq\emptyset$. (O1) is Lemma~\ref{bpP:p11a:finite-divisor}.

\proofstep{Hypothesis (C2)}\label{stepP:p11a:il-c2} Apply Proposition~\ref{bpP:p11a:canonical-step} to the stage $\nu$ with $E=K_{\nu+1}$, $t=t_{\nu+1}$, $B=E+t=B_{\nu+1}$ ($\subset B_\nu$; the restrictions of $h_\nu,\omega_\nu,R_\nu$ to $B_{\nu+1}$), $r=r_\nu$, $s=s_\nu$, $T=T_\nu$, $d=d_\nu$, $\delta=\delta_\nu$, $\sigma=\sigma_\nu$, $q$, $D=D_\nu$, and with $\varepsilon:=\varepsilon_\nu$ (the proposition holds verbatim for any upper bound $\varepsilon\ge\nrmm R{}$: its proof uses $\lVert R\rVert\le\varepsilon$ only, through Proposition~\ref{bpP:p11a:cutoff}; here $\varepsilon_\nu>0$, so $A_1>0$ as Lemma~\ref{bpP:p11a:affine-flow} requires). The divisor part of (C2) is (O1). The bound on $D\omega_\nu$ on $B_{\nu+1}$: $L=2C_1$ by $(\mathrm P_\nu)$(a), and $t_{\nu+1}T_\nu L\le D_\nu/2$ is \ref{bpP:p11a:SC5}.

\proofstep{Hypotheses (C1) and the loss}\label{stepP:p11a:il-c1} $r_{\nu+1}=r_\nu-d_\nu-\delta_\nu-2\sigma_\nu\ge s/2>0$ and $d_\nu\le\frac14$, $E_\nu\le\alpha\le1$ (Lemma~\ref{bpP:p11a:strip-budget}); $q\le\frac12$ (Lemma~\ref{bpP:p11a:schedule-ratios}); $\delta_\nu\le s/24\le1$.

\proofstep{Hypothesis (C3) and the recursion}\label{stepP:p11a:il-c3} By \Kref{bp:p1:exp-sums}(b) (as in Section~\ref{bpP:p11a:sec-hom}), $\LamD\delta\le(1+\frac1{\pi\delta})^N-1\le(2/\delta)^N$ for $0<\delta\le1$ (as $1+\frac1\pi<2$). Hence the relative coefficient $Q_\nu:=C_Q\LamD{\delta_\nu}\varepsilon_\nu/(\pi D_\nu s_\nu\sigma_\nu)\le Q^\sharp_\nu=Q_0g_1^\nu$ (\ref{bpP:p11a:SC4}). Since $\varepsilon_\nu=\varepsilon_0\varpi^\nu$ is a prescribed sequence (the induction below is that of Lemma~\ref{bpP:p11a:geom-recursion}, with $c_0=2q^2+2\alpha=\frac38\varpi$, $P_0=Q_0/\varepsilon_0$, $p=g_1/\varpi$, $p\varpi=g_1\le1$), only the inequality $Q_\nu\le\varpi/2$ is needed: it is $Q_\nu\le Q_0g_1^\nu\le Q_0\le\varpi/2$ (\ref{bpP:p11a:A4}). Since $C_Q\ge8(N+1)\ge4(N+1)(1+E_\nu)$, $Q_\nu\le\varpi/2<1$ gives $4(N+1)\LamD{\delta_\nu}(1+E_\nu)\varepsilon_\nu\le\pi D_\nu s_\nu\sigma_\nu$: (C3).

\proofstep{The new remainder}\label{stepP:p11a:il-rem} Proposition~\ref{bpP:p11a:canonical-step}(iii), second form: $\nrmm{R_+}{}\le\varepsilon_\nu(2q^2+2E_\nu)+\frac{18N(4N+5)\LamD{\delta_\nu}}{\pi D_\nu s_\nu\sigma_\nu}\varepsilon_\nu^2=\varepsilon_\nu(2q^2+2E_\nu+Q_\nu)\le\varepsilon_\nu(\frac18\varpi+\frac14\varpi+\frac12\varpi)\le\varepsilon_\nu\varpi=\varepsilon_{\nu+1}$, by $2q^2=\varpi/8$, $2E_\nu\le2\alpha=\varpi/4$ (Lemma~\ref{bpP:p11a:schedule-ratios}). The domain of $R_+$ is $\Ddom{r_{\nu+1}}{qs_\nu}{B_{\nu+1}}=\Ddom{r_{\nu+1}}{s_{\nu+1}}{B_{\nu+1}}$; periodicity and reality are in the proposition. This is $(\mathrm P_{\nu+1})$(b), and (O5) with (c) by composition.

\proofstep{The increments on the whole collar}\label{stepP:p11a:il-inc} Let $R_*=\cutT_{T_\nu}(R_\nu^{(\le1)})$ be formed on $\Ddom{r_\nu}{s_\nu}{B_\nu}$ (Proposition~\ref{bpP:p11a:cutoff} with $B=B_\nu$; it does not involve the divisor), $\Delta h_\nu=\langle R_{*0}\rangle$, $\Delta\omega_\nu=\langle R_{*1}\rangle$. These are holomorphic on $B_\nu$ (Step~\ref{stepP:p11a:cu-coeff}), real at real points (Proposition~\ref{bpP:p11a:cutoff}(c)), $\lvert\Delta h_\nu\rvert\le(1+E_\nu)\varepsilon_\nu\le2\varepsilon_\nu$ and $\lvert\Delta\omega_\nu\rvert_\infty\le(1+E_\nu)\varepsilon_\nu/s_\nu\le\ell_\nu$ (Proposition~\ref{bpP:p11a:cutoff}(a) and $k=0$ of \Kref{bp:p1:decay}); on $B_{\nu+1}$ they are the $\Delta h,\Delta\omega$ of Proposition~\ref{bpP:p11a:canonical-step}(i), which gives (O2).

\proofstep{$D\omega_{\nu+1}$}\label{stepP:p11a:il-dom} For $b\in B_{\nu+1}$ and $j\le\nu$ pick $b'\in K_{\nu+1}\subset K_j$ with $\lvert b-b'\rvert<t_{\nu+1}\le t_j/2$: $b\in K_j+\frac{t_j}2$. Lemma~\ref{bpP:p11a:param-cauchy} ($E=K_j$, $t=t_j$, $t'=t_j/2$, $u=1$): $\lVert D\Delta\omega_j(b)\rVert\le\ell_j(t_j/2)^{-1}=4\varepsilon_j/(s_jt_j)$. So $\lVert D\omega_{\nu+1}\rVert\le C_1+\sum_{j\le\nu}4\varepsilon_j/(s_jt_j)\le2C_1$ by \ref{bpP:p11a:SC3}: $(\mathrm P_{\nu+1})$(a). ($\Delta h_\nu$ is holomorphic on $B_\nu\supset B_{\nu+1}$ and real.)

\proofstep{The map, (O3), (O4)}\label{stepP:p11a:il-map} Proposition~\ref{bpP:p11a:canonical-step}(ii),(iv),(v) with the displayed $A_1=A_{1,\nu}$, $A_\eta=A_{\eta,\nu}$, $\kappa=\kappa_\nu\le\frac12$ ((C3) gives $4(N+1)A_1\le\sigma$, i.e.\ $\kappa\le\frac12$). The unscaled derivative bound: in the block form of Lemma~\ref{bpP:p11a:affine-flow}(c),(d) the first row has norm $\lvert D\Theta-\mathbf 1\rvert\le2\kappa$ and the second row $\lvert\partial_\theta(\mathbf M\eta+m)\rvert+\lvert\mathbf M-\mathbf 1\rvert\le\kappa s_\nu/\sigma_\nu+2\kappa$, maximal row sum (review P11A2-04). \end{proof}

%%%%%%%%%%%%%%%%%%%%%%%%%%%%%%%%%%%%%%%%%%%%%%%%%%%%%%%%%%%%%%%%%%%%%%%%%%%%%%%%
\subsection{Parameter selection}\label{bpP:p11a:sec-select}

\begin{theorem}[Parameter selection; Rus01 Thm 16.7 p.~179, proof \S20.7 pp.~201--203]\label{bpP:p11a:parameter-selection}
Let Assumption~\ref{bpP:p11a:sel-data} hold and $0\le\varepsilon\le\varepsilon_{\mathrm{adm}}$ (Definition~\ref{bpP:p11a:adm-params}). Then the iteration of Proposition~\ref{bpP:p11a:inductive-lemma} can be carried out for all $\nu\ge0$, and:
\begin{enumerate}[label=(\alph*)]
\item $K=K_0\supset K_1\supset\cdots$, $K_\infty:=\bigcap_\nu K_\nu$ is compact; for $b\in K_\infty$ the limits $\omega_*(b)=\omega(b)+\sum_{j\ge0}\Delta\omega_j(b)$ and $h_*(b)=\sum_{j\ge0}\Delta h_j(b)$ exist, $\lvert\omega_*(b)-\omega_m(b)\rvert_\infty\le S_m\le\gamma T_m^{-\tau-1}$, $\lvert h_*(b)-h_m(b)\rvert\le4\varepsilon_m$, and $\omega_*:K_\infty\to\R^N$, $h_*:K_\infty\to\R$ are continuous;
\item the set $H_\infty:=\bigcap_{\nu\ge0}H^{(\nu)}=\{b\in K:\lvert k\cdot\tilde\omega_{\nu(k)}(b)\rvert\ge2\gamma\lvert k\rvert_1^{-\tau}\ \forall k\ne0\}$ is compact, $H_\infty\subset K_\infty$, and $\operatorname{meas}H_\infty\ge\operatorname{meas}K-\varepsilon^*$;
\item for every $b\in H_\infty$ the vector $\omega_*(b)$ is $(\gamma,\tau)$-Diophantine in the frozen sense: $\lvert k\cdot\omega_*(b)\rvert\ge\gamma\lvert k\rvert_1^{-\tau}$ for all $k\in\Z^N\setminus\{0\}$.
\end{enumerate}
In particular $\operatorname{meas}K_\infty\ge\operatorname{meas}H_\infty\ge\operatorname{meas}K-\varepsilon^*$, and the quantifier order is data $\to\varepsilon^*\to(\gamma,\varepsilon_{\mathrm{adm}})\to f$: $\gamma,\varepsilon_{\mathrm{adm}}$ depend on the data and $\varepsilon^*$ only. No premise equals a conclusion: the Diophantine property is an output.
\source{Rus01 Thm 16.7 (p.~179: $K_\infty\ne\emptyset$, $\operatorname{meas}K_\infty\ge\operatorname{meas}H\ge\operatorname{meas}K-\varepsilon^*$) and Thm 20.5 (pp.~201--203); the explicit constants (20.33)--(20.35) are replaced by Definition~\ref{bpP:p11a:adm-params} (deviation D-9); mode-wise selection with the finite-stage invariant (Remark~\ref{bpP:p11a:i3-corrected}), after Remark~\ref{bpP:p11a:consumers}.}
\end{theorem}
\begin{proof}
\proofstep{Infinite iteration}\label{stepP:p11a:ps-iter} Proposition~\ref{bpP:p11a:inductive-lemma} gives $(\mathrm P_0)$ and $(\mathrm P_\nu)\Rightarrow(\mathrm P_{\nu+1})$, hence $(\mathrm P_\nu)$ for all $\nu$ with nonempty compact $K_\nu$ decreasing. $K_\infty$ is a decreasing intersection of nonempty compact sets (so nonempty), closed and bounded.

\proofstep{Limits}\label{stepP:p11a:ps-lim} For $b\in K_\infty$: $b\in K_j\subset B_j$ for all $j$ and $\sum_j\lvert\Delta\omega_j(b)\rvert_\infty\le\sum\ell_j=S_0<\infty$ (\ref{bpP:p11a:SC2}); $\sum_j\lvert\Delta h_j\rvert\le2\sum\varepsilon_j<\infty$. The partial sums $\omega_\nu$ converge uniformly on $K_\infty$ (each $\Delta\omega_j$ is continuous on $K_\infty\subset B_j$, uniform Weierstrass test, \Kref{bp:p1:uniform-limits}) so $\omega_*$ is continuous; the tail bound is $\sum_{j\ge m}\ell_j=S_m$, and $\sum_{j\ge m}2\varepsilon_j\le4\varepsilon_m$ ($\varpi\le\frac12$).

\proofstep{The set $H_\infty$}\label{stepP:p11a:ps-H} $H^{(\nu)}$ is compact and $H^{(\nu)}\supset H^{(\nu+1)}$ (more modes are tested; $\tilde\omega_{\nu(k)}$ is the same map for a given $k$ whichever $\nu$ is used, because the chain extension $\tilde\omega_j$ depends only on the first $j$ stages). So $H_\infty=\bigcap H^{(\nu)}$ is compact; by \ref{bpP:p11a:y-down} and Lemma~\ref{bpP:p11a:finite-measure}, $\operatorname{meas}H_\infty=\lim\operatorname{meas}H^{(\nu)}\ge\operatorname{meas}K-\varepsilon^*$. Also $H_\infty\subset K_\nu$ for every $\nu$ (Lemma~\ref{bpP:p11a:finite-measure}, $H^{(\nu)}\subset K_{\nu+1}$), so $H_\infty\subset K_\infty$.

\proofstep{Diophantine property}\label{stepP:p11a:ps-dio} Let $b\in H_\infty$, $k\ne0$, $m=\nu(k)$. Then $u:=\omega_m(b)=\tilde\omega_m(b)$ ($b\in K_m$) satisfies $\lvert k\cdot u\rvert\ge2\gamma\lvert k\rvert_1^{-\tau}$ ($b\in H^{(m)}$), and $v:=\omega_*(b)$ satisfies $\lvert u-v\rvert_\infty\le S_m\le\gamma T_m^{-\tau-1}\le\gamma\lvert k\rvert_1^{-\tau-1}$ (part (a), $\lvert k\rvert_1\le T_m$). Lemma~\ref{bpP:p11a:limit-bridge} gives $\lvert k\cdot\omega_*(b)\rvert\ge\gamma\lvert k\rvert_1^{-\tau}$.
\end{proof}

%%%%%%%%%%%%%%%%%%%%%%%%%%%%%%%%%%%%%%%%%%%%%%%%%%%%%%%%%%%%%%%%%%%%%%%%%%%%%%%%
\subsection{What File 11a-4 takes: summable budgets}\label{bpP:p11a:sec-export}

\begin{lemma}[Exported sums]\label{bpP:p11a:exported-sums}
In the situation of Theorem~\ref{bpP:p11a:parameter-selection} (so \ref{bpP:p11a:A1}--\ref{bpP:p11a:A7}, with $Q_0$ as in \ref{bpP:p11a:SC4}) let $A_{1,\nu},A_{\eta,\nu},\kappa_\nu,\Lambda_{\mathbb R,\nu}$ be as in Proposition~\ref{bpP:p11a:inductive-lemma}\ref{bpP:p11a:O3}. Then for all $\nu\ge0$:
\begin{enumerate}[label=(E\arabic*),ref=(E\arabic*)]
\item\label{bpP:p11a:E1} $A_{1,\nu}\le2\sigma_\nu Q_\nu/C_Q$, $\kappa_\nu\le Q_\nu/16$, $A_{\eta,\nu}=\frac12\kappa_\nu s_\nu$, $\Lambda_{\mathbb R,\nu}\le\dfrac{4\pi}{C_Q}T_0\varsigma_0Q_\nu$, where $Q_\nu\le Q_0g_1^\nu\le Q_02^{-5\nu}$;
\item\label{bpP:p11a:E2} (angle, action and Lipschitz budgets, unscaled) $\sum_\nu(2\kappa_\nu+\kappa_\nu s_\nu/\sigma_\nu)\le2^{-12}$, $\ \sum_\nu\Lambda_{\mathbb R,\nu}\le2^{-12}$, $\ \sum_\nu A_{1,\nu}\le3\varsigma_0Q_0/C_Q$, $\ \sum_\nu A_{\eta,\nu}\le s_0Q_0/16$;
\item\label{bpP:p11a:E3} (tails against the collars) $\displaystyle\sup_{\nu\ge0}\ t_{\nu+1}^{-1}\sum_{j\ge\nu}\bigl(A_{1,j}+A_{\eta,j}\bigr)\le2^{-10}$, and $\sum_{j\ge\nu}\Lambda_{\mathbb R,j}\le2^{-10}$;
\item\label{bpP:p11a:E4} (parameter Cauchy) on $K_{\nu+1}+t_{\nu+1}/2$: $\lVert D_b(X_\nu-\mathrm{id})\rVert\le2(A_{1,\nu}+A_{\eta,\nu})/t_{\nu+1}$.
\end{enumerate}
\source{Ours; replaces Rus01 15.1--15.2 (the summability $\sum C_\nu T_\nu<\infty$, (14.20)) and (A3) item 3 of the review; the unscaled form is the correction P11A2-04 of the review of File~11a-2.}
\end{lemma}
\begin{proof}
\proofstep{(E1)}\label{stepP:p11a:es-1} $A_{1,\nu}=(1+E_\nu)\sigma_\nu Q_\nu/C_Q\le2\sigma_\nu Q_\nu/C_Q$ by definition of $Q_\nu$. $\kappa_\nu=2(N+1)A_{1,\nu}/\sigma_\nu\le4(N+1)Q_\nu/C_Q\le Q_\nu/16$ since $C_Q=18N(4N+5)\ge64(N+1)$. $A_{\eta,\nu}=(N+1)s_\nu A_{1,\nu}/\sigma_\nu=\frac12\kappa_\nu s_\nu$. $\Lambda_{\mathbb R,\nu}\le4T_\nu\varepsilon_\nu\LamD{\delta_\nu}/(D_\nu s_\nu)=\frac{4\pi}{C_Q}T_\nu\sigma_\nu Q_\nu$, because $r_\nu-d_\nu\ge\delta_\nu+s/2>\delta_\nu$ (Lemma~\ref{bpP:p11a:strip-budget}) and $\LamD\cdot$ decreases; and $T_\nu\sigma_\nu=T_0\varsigma_0(\Theta\zeta)^\nu=T_0\varsigma_0$. $Q_\nu\le Q^\sharp_\nu=Q_0g_1^\nu$ (Step~\ref{stepP:p11a:il-c3}).

\proofstep{(E2)}\label{stepP:p11a:es-2} $\sum_\nu Q_\nu\le\frac{32}{31}Q_0$. $s_\nu/\sigma_\nu=(s_0/\varsigma_0)(2q)^\nu\le s_0/\varsigma_0\le t_0/\varsigma_0$ (\ref{bpP:p11a:A3}), so $\sum(2\kappa_\nu+\kappa_\nu s_\nu/\sigma_\nu)\le\frac1{16}\cdot\frac{32}{31}Q_0(2+t_0/\varsigma_0)\le Q_0\Xi/7\le2^{-12}$ (\ref{bpP:p11a:A4}, $2+t_0/\varsigma_0\le2\Xi$). $\sum\Lambda_{\mathbb R,\nu}\le\frac{4\pi}{C_Q}T_0\varsigma_0\frac{32}{31}Q_0\le\frac{13}{162}Q_0\Xi\le2^{-12}$. $\sum A_{1,\nu}\le\frac{2\varsigma_0Q_0}{C_Q}\sum(\zeta g_1)^\nu\le3\varsigma_0Q_0/C_Q$; $\sum A_{\eta,\nu}\le\frac1{32}s_0Q_0\sum(qg_1)^\nu\le s_0Q_0/16$.

\proofstep{(E3)}\label{stepP:p11a:es-3} $\sum_{j\ge\nu}A_{1,j}\le\frac{2\varsigma_0Q_0}{C_Q}\sum_{j\ge\nu}(\zeta g_1)^j\le\frac{3\varsigma_0Q_0}{C_Q}(\zeta g_1)^\nu$ and $t_{\nu+1}=t_02^{-(\tau+1)}(2^{-(\tau+1)})^\nu$, so the quotient is $\le\frac{3\varsigma_0Q_0\Theta^{\tau+1}}{C_Qt_0}g_4^\nu\le\frac{3}{C_Q}Q_0\Xi\le2^{-10}$ ($g_4\le1$). Similarly $\sum_{j\ge\nu}A_{\eta,j}\le\frac{s_0Q_0}{16}(qg_1)^\nu$ and $s_0\le t_0$, $qg_1\le g_4$: quotient $\le Q_0\Xi/16$. Sum of the two is $\le2^{-10}$. $\sum_{j\ge\nu}\Lambda_{\mathbb R,j}\le2^{-12}$ by (E2).

\proofstep{(E4)}\label{stepP:p11a:es-4} Proposition~\ref{bpP:p11a:canonical-step}(v) with $t'=\frac12t_{\nu+1}$, $j=1$ gives $\le A/(t_{\nu+1}/2)$ for each component, with $A=A_{1,\nu},A_{\eta,\nu}$; the max norm of the pair is at most the sum.
\end{proof}

\paragraph{What File 11a-4 must take from this file.}
\begin{enumerate}[label=(T\arabic*)]
\item The sets $K_\nu\supset H_\infty$ and the maps $X_\nu$, holomorphic on $\Ddom{r_{\nu+1}}{s_{\nu+1}}{B_{\nu+1}}$, with $r_{\nu+1}\ge s/2$, $s_{\nu+1}=s_0q^{\nu+1}$, $B_{\nu+1}=K_{\nu+1}+t_{\nu+1}$; the composition $X_0\circ\dots\circ X_\nu$ is defined on $\Ddom{r_{\nu+1}}{s_{\nu+1}}{B_{\nu+1}}$ with values in $\Ddom s{s_0}{B_0}$ (domain bookkeeping of Proposition~\ref{bpP:p11a:inductive-lemma}\ref{bpP:p11a:O3}). The remainders satisfy $\lvert R_\nu\rvert\le\varepsilon_0\varpi^\nu$ on $\Ddom{r_\nu}{s_\nu}{B_\nu}$ and $H_\nu=h_\nu+\omega_\nu\cdot\eta+R_\nu$. The shrinking-domain estimates give exactly two \emph{zero-section limits}: for $b\in K_\infty$ and every $\theta$ in the strip, $\lvert R_\nu(\theta,0,b)\rvert\le\varepsilon_\nu\to0$, and, by Cauchy in $\eta$ on $\lvert\eta\rvert<s_\nu/2$, $\lvert\partial_\eta R_\nu(\theta,0,b)\rvert\le2\varepsilon_\nu/s_\nu=\ell_\nu\to0$. Hence the limit torus (the image of the zero section under $X_0\circ\dots\circ X_\nu$), on which $H_\nu=h_\nu+\omega_\nu\cdot\eta+R_\nu$ has vanishing remainder and vanishing first $\eta$-jet of the remainder in the limit, is what this file delivers. It does \emph{not} deliver convergence of $\partial_\eta^2R_\nu$, nor a limiting normal form $h_*+\omega_*\cdot\eta+O(\eta^2)$ on a fixed action neighbourhood: $s_\nu\to0$ and $\varepsilon_\nu/s_\nu^2=(\varepsilon_0/s_0^2)\,16^\nu\to\infty$. If File~11a-4 wants an off-section normal form, it must prove the convergence and extension of the affine action coefficients $\Delta h_\nu,\Delta\omega_\nu$ (and of the maps) on a common domain; the physical consumer needs only the torus and its first-jet identity.
\item Displacement of the defect: $\sum_\nu A_{1,\nu}\le3\varsigma_0Q_0/C_Q$ (angles), $\sum_\nu A_{\eta,\nu}\le s_0Q_0/16$ (actions), both explicit in the data (E2); the actions $I=b+\eta$ differ from $b$ by the action defect, whose total size is at most $s_0$ (the initial translation) plus $\sum A_{\eta,\nu}$ on $\lvert\eta\rvert<s_0$; for $\eta=0$ it is $\le\sum A_{\eta,\nu}$.
\item The \emph{unscaled} derivative budgets, not the stage-norm bound $2\kappa_\nu$: $\lvert DX_\nu-\mathbf 1\rvert_\infty\le2\kappa_\nu+\kappa_\nu s_\nu/\sigma_\nu$ in $(\theta,\eta)$ with $\sum\le2^{-12}$ (E2), and $\lvert D\Theta_\nu-\mathbf 1\rvert_\infty\le e^{\Lambda_{\mathbb R,\nu}}-1$ on the reals with $\sum\Lambda_{\mathbb R,\nu}\le2^{-12}$; these are the inputs for the $\lambda$-Lipschitz bound in $\theta$ (the compositions of maps close to the identity in the fixed sup-norm). The parameter derivatives (E4) and tails (E3): $\sum_{j\ge\nu}(A_{1,j}+A_{\eta,j})\le2^{-10}t_{\nu+1}$ (the plan's ``$\kappa$-versus-$\tau$'' condition: it holds because $g_4\le2^{-5}$ and $Q_0\Xi\le2^{-10}$). The local-to-global argument takes $b,b'\in K_\infty$ with $\lvert b-b'\rvert_\infty<t_{\nu+1}/2$: the segment $[b,b']$ lies in $K_{\nu+1}+t_{\nu+1}/2$ (within $t_{\nu+1}/2$ of $b\in K_{\nu+1}$), where (E4) applies; no convexity of $K_\infty$ is used.
\item The frequency: $\omega_*$ on $K_\infty$ and $w_b=\omega_*(b)$ is $(\gamma,\tau)$-Diophantine exactly on a set containing $H_\infty$ (Theorem~\ref{bpP:p11a:parameter-selection}(c)); $\lvert\omega_*-\omega_\nu\rvert_\infty\le S_\nu$ on $K_\infty$.
\item The inequalities \ref{bpP:p11a:A1}--\ref{bpP:p11a:A7}, \ref{bpP:p11a:E2}--\ref{bpP:p11a:E3} are \emph{part of the admissibility} $\varepsilon\le\varepsilon_{\mathrm{adm}}$ (the factor $\Xi$ in $\lambda_{\max}$ is what makes \ref{bpP:p11a:E2}, \ref{bpP:p11a:E3} hold); 11a-4 adds no further smallness unless its $\lambda$-Lipschitz constant must be $<2^{-10}$, in which case it multiplies $\lambda_{\max}$ by a constant.
\end{enumerate}

%%%%%%%%%%%%%%%%%%%%%%%%%%%%%%%%%%%%%%%%%%%%%%%%%%%%%%%%%%%%%%%%%%%%%%%%%%%%%%%%
\subsection{The polynomial budget and its numerical check (risk R1)}\label{bpP:p11a:sec-num}

Part~11 chooses the order of its finite preparation after the exponent of the budget, so the budget must be a proved lower bound in the exact form in which Part~11 uses it: every datum except one is held fixed and the \emph{amount} $\beta$ tends to zero. The only data on which $\varepsilon_{\mathrm{adm}}$ depends are $(N,\mu,\rho,\tau)$ (fixed integers and exponents) and $(\vartheta,s,\beta,\varepsilon^*,M,M_0,C_1,d)$, where $d=\operatorname{diam}_\infty K$ and $d,M$ enter only through $\gamma$ ($\beta$ enters through $\gamma$ and directly through $\theta_{\max}$). The next proposition bounds $\varepsilon_{\mathrm{adm}}$ from below along an arbitrary one-parameter family of such data, in terms of explicit polynomial rates.

\begin{proposition}[Polynomial lower bound for $\varepsilon_{\mathrm{adm}}$]\label{bpP:p11a:adm-budget}
Fix $N,\mu\ge1$, $\rho\in\{1,2\}$ and $\tau>\frac{N\mu}\rho-1$. Let $s_\tau:=\frac{\rho(\tau+1)}\mu-N>0$ (the lattice gap of Corollary~\ref{bpP:p11a:power-law}; it is a constant, not the strip width $s$). Let $c\in(0,1]$ and exponents $a_\vartheta,a_s,a_\beta,a_*,a_M,a_0,a_C,a_d\ge0$ be given, and let, for every $u\in(0,1]$, data $(\vartheta,s,\beta,\varepsilon^*,M,M_0,C_1,d,K,\ldots)$ as in Assumption~\ref{bpP:p11a:sel-data} (with $C_1\ge1$) be given such that
\begin{gather*}\vartheta\ge c\,u^{a_\vartheta},\quad s\ge c\,u^{a_s},\quad \beta\ge c\,u^{a_\beta},\quad\varepsilon^*\ge c\,u^{a_*},\\ M\le c^{-1}u^{-a_M},\quad M_0\le c^{-1}u^{-a_0},\quad C_1\le c^{-1}u^{-a_C},\quad d\le c^{-1}u^{-a_d}.\end{gather*}
Put
\[D=Na_d+a_M+a_\vartheta+\bigl(1+\tfrac1\mu\bigr)a_\beta,\qquad G=\max\Bigl\{\frac{a_\beta}\rho,\ \frac\mu\rho\,(a_*+D)\Bigr\},\qquad L=\begin{cases}a_s+1&a_s>0,\\0&a_s=0,\end{cases}\]
\begin{gather*}J=\max\{a_\vartheta,\ G+(\tau+1)L+a_C\},\qquad X=\max\{J,L\},\\ H=\max\bigl\{J,\ G+(N+1)a_s+\tau L+X,\ a_\beta+a_0+\mu J\bigr\},\qquad B=2H+a_C.\end{gather*}
Then there is $c'>0$, depending only on $(N,\mu,\rho,\tau,c)$ and on the strip rate $a_s$ (for $a_s=0$: on $(N,\mu,\rho,\tau,c)$ alone; in no case on $u$, on the family once the rates are fixed, or on the iteration index), with
\[\varepsilon_{\mathrm{adm}}\ \ge\ c'\,u^{B}\qquad(u\in(0,1]).\]
\source{Ours; this replaces the monomial asymptotics asserted in the first version of this section (reviews P11A3-01, -02). The exponent formula is the review's, with one more entry checked: the factor $1+1/s_\tau$ in $C_{\mathrm{Dio}}$ is a constant, not the strip width.}
\end{proposition}
\begin{proof}
Below $C,c_i>0$ denote constants depending only on $(N,\mu,\rho,\tau,c)$ and $a_s$ (this is needed only in Step~\ref{stepP:p11a:ab-T0}; all other steps are free of the rates); every exponent is $\ge0$ and $u\le1$, so enlarging an exponent weakens a bound of the form $u^{-e}$ or $u^{e}$ in the right direction.

\proofstep{The loss $\gamma$}\label{stepP:p11a:ab-gamma} Since $M/\beta\le c^{-2}u^{-a_M-a_\beta}$, $d\le c^{-1}u^{-a_d}$ and $\vartheta\ge cu^{a_\vartheta}$: $C_{\mathrm{sub}}=2^{N+2}\mu(\sqrt Nd)^{N-1}\frac M\beta(1+(4+\vartheta^{-1})\sqrt Nd)\le Cu^{-(N-1)a_d-a_M-a_\beta-a_d-a_\vartheta}$ (the last bracket is $\le C(1+u^{-a_\vartheta})(1+u^{-a_d})\le Cu^{-a_\vartheta-a_d}$). Hence $C_{\mathrm{Dio}}=C_{\mathrm{sub}}\beta^{-1/\mu}N^{\rho/2\mu}2\cdot3^{N-1}(1+1/s_\tau)\le Cu^{-D}$, because $\beta^{-1/\mu}\le c^{-1/\mu}u^{-a_\beta/\mu}$. So both entries of $\gamma_*(\varepsilon^*)=\min\{\frac12N^{-1/2}(\beta/4)^{1/\rho},(\varepsilon^*/C_{\mathrm{Dio}})^{\mu/\rho}\}$ are $\ge c_1u^{a_\beta/\rho}$, respectively $\ge c_1u^{(\mu/\rho)(a_*+D)}$, and $\gamma=\gamma_*/2\ge c_2u^{G}$.

\proofstep{The cutoff}\label{stepP:p11a:ab-T0} $T_0=\max\{1,\frac{16}{\pi s}(A_0+N\log\frac{32N}s)\}$ with $A_0$ depending on $(N,\mu,\tau)$ only. If $a_s=0$, $s\ge c$ and $T_0\le C$. If $a_s>0$, $\log\frac{32N}s\le\log\frac{32N}c+a_s\log\frac1u\le C_0(1+a_s)u^{-1}$ (as $\log\frac1u\le\frac1u$) and $\frac1s\le c^{-1}u^{-a_s}$, so $T_0\le C_0(1+a_s)u^{-a_s-1}$; the factor $1+a_s$ is absorbed in $C$, which is why $C$ and hence $c'$ may depend on $a_s$. In both cases $T_0\le Cu^{-L}$, and $a=\gamma T_0^{-\tau-1}\ge c_3u^{G+(\tau+1)L}$.

\proofstep{The collar and $\Xi$}\label{stepP:p11a:ab-t0} $t_0=\min\{\vartheta/\sqrt N,a/(4C_1)\}\ge c_4u^{J}$ (use $C_1\le c^{-1}u^{-a_C}$). Since $\varsigma_0=s/24\le1$, $t_0\le1$ and $\varsigma_0\ge c_5u^{a_s}$: $\Xi=2^{\tau+1}(1+\frac{t_0}{\varsigma_0}+\frac{\varsigma_0}{t_0}+T_0\varsigma_0)\le C(1+u^{-a_s}+u^{-J}+u^{-L})\le Cu^{-X}$ ($a_s\le L$).

\proofstep{The three terms}\label{stepP:p11a:ab-lam} (1) $C_1t_0/8\ge c_4u^J/8$ ($C_1\ge1$). (2) $\min\{\varpi/2,2^{-10}/\Xi\}\ge c_6u^{X}$ and $\frac{\pi\gamma\varsigma_0^{N+1}}{C_Q2^NT_0^\tau}\ge c_7u^{G+(N+1)a_s+\tau L}$, so the second term is $\ge c_8u^{G+(N+1)a_s+\tau L+X}$. (3) For $\rho=1$, $\theta_{\max}=\beta/2\ge\frac c2u^{a_\beta}$. For $\rho=2$: if $\beta\le2M_0$ then $2M_0+\beta\le4M_0$ and $\beta/(2^{\mu+1}(2M_0+\beta))\ge2^{-\mu-3}\beta/M_0$; otherwise $2M_0+\beta<2\beta$ and it is $\ge2^{-\mu-2}$. Hence $\theta_{\max}\ge2^{-\mu-3}\min\{\beta,\beta/M_0,1\}\ge c_9u^{a_\beta+a_0}$. In both cases the third term is $\ge c_{10}u^{a_\beta+a_0+\mu J}$ by $t_0\ge c_4u^J$. Altogether $\lambda_{\max}\ge c_{11}u^{H}$.

\proofstep{Conclusion}\label{stepP:p11a:ab-concl} $\varepsilon_{\mathrm{adm}}=\lambda_{\max}^2/(2NC_1)\ge c_{11}^2u^{2H}\,c\,u^{a_C}/(2N)$.
\end{proof}

\begin{remark}[Why the hypotheses are all needed]\label{bpP:p11a:adm-budget-necessary}
The dependence of $c'$ on $a_s$ cannot be removed: take fixed $(N,\mu,\rho,\tau,c)$ and the family with only $s=c\,u^{A}$ varying ($a_s=A$, all other rates $0$). Specialize to $(N,\mu,\rho,\tau,c)=(2,3,2,3,\frac1{10})$, $h(I)=\frac12\lvert I\rvert_2^2$, $K=[0,\frac12]^2$, $\vartheta=\beta=\varepsilon^*=\frac1{10}$, $M=4$, $M_0=C_1=1$ and $s(u)=u^A/10$ (for other indices the displayed formula for $B$ need not hold: in general $B=2\max\{(N+2\tau+2)A+2\tau+1,\ \mu(\tau+1)(A+1)\}$ and the third entry need not dominate). Then $B=24(A+1)$, and at the same $u=e^{-1}$ one has $T_0\ge c_0e^{A}(A+b)$ for fixed $c_0,b>0$, so $\varepsilon_{\mathrm{adm}}\le\varepsilon^{(3)}\le CT_0^{-2\mu(\tau+1)}$ gives $\varepsilon_{\mathrm{adm}}/u^{B}\le C'(A+b)^{-24}\to0$ as $A\to\infty$ (review P11A3R2-01). The corollary below has $a_s=0$. Dropping the upper bound on $M$ destroys the bound: take $N=2$, $h(I)=\frac12\lvert I\rvert_2^2$ on a fixed box with $\mu=3$, $\rho=2$, fixed $\vartheta,s,\beta,\varepsilon^*$, and let the admissible upper bound $M=M(u)$ for the derivatives of $g_\rho[\omega]$ be $e^{1/u}$ (any larger bound is again a bound); then $\gamma\propto M^{-3/2}$ and, as the closeness branch eventually binds, $\varepsilon_{\mathrm{adm}}\propto M^{-9}$, which is smaller than every $cu^{B}$. Dropping the normalisation $C_1\ge1$ does the same: for $h(I)=I$ ($N=1$) every $C_1>0$ is a valid Hessian bound, and $\varepsilon^{(1)}=C_1t_0^2/(128N)$ is proportional to $C_1$ for small $C_1$. The diameter $d$ and the size $M_0$ of $\omega$ enter likewise and are therefore data of the family.
\end{remark}

\begin{corollary}[The budget in the form used by Part~11]\label{bpP:p11a:adm-budget-beta}
Fix $N,\mu\ge1$, $\rho\in\{1,2\}$, $\tau>\frac{N\mu}\rho-1$, and fix \emph{all} data of Assumption~\ref{bpP:p11a:sel-data} --- $K$, $\vartheta$, $s$, $\varepsilon^*$, $M$, $M_0$, $C_1$ --- except the amount $\beta$, which is free in $(0,\beta_{\max}]$, $\beta_{\max}=\min\{1,M\}$. Put
\[G_0=\frac{\mu+1}\rho,\qquad B_0=2\max\{2G_0,\ 1+\mu G_0\}.\]
Then there is $c_*>0$, depending only on the fixed data, with $\varepsilon_{\mathrm{adm}}\ge c_*\,\beta^{B_0}$ for all $\beta\in(0,\beta_{\max}]$. For $\rho=2$, $B_0=\mu(\mu+1)+2$.
\source{Ours; the form of the polynomial budget (K5) of Definition~\ref{bpP:p11:kam-input}: the fixed-data case $a_\beta=1$, all other exponents $0$, of Proposition~\ref{bpP:p11a:adm-budget}.}
\end{corollary}
\begin{proof}
\proofstep{Specialise}\label{stepP:p11a:abb-spec} Take $u=\beta\in(0,1]$, $a_\beta=1$, $a_\vartheta=a_s=a_*=a_M=a_0=a_C=a_d=0$ and $c$ the minimum of $1$, $\vartheta$, $s$, $\varepsilon^*$, $M^{-1}$, $M_0^{-1}$, $C_1^{-1}$, $d^{-1}$ (all fixed). Then $D=1+\frac1\mu$, $G=\max\{\frac1\rho,\frac\mu\rho(1+\frac1\mu)\}=\frac{\mu+1}\rho=G_0$, $L=0$, $J=\max\{0,G_0\}=G_0$, $X=G_0$, $H=\max\{G_0,\ 2G_0,\ 1+\mu G_0\}$, $B=2H$. Since $2G_0\ge G_0$ this is $B_0$. For $\rho=2$, $G_0=\frac{\mu+1}2$ and $1+\mu G_0\ge2G_0$ iff $\mu\ge1$, so $B_0=2+\mu(\mu+1)$.
\end{proof}

\paragraph{Reading of the exponent.} Each of the three branches of Remark~\ref{bpP:p11a:adm-branches} contributes one entry of $H$: the first and third through $J$ (the collar $t_0\gtrsim u^{J}$), the second through $G+(N+1)a_s+\tau L+X$ (the quadratic branch, including $\Xi$). The schedule contributes only constants, powers of $2$ ($\varpi=2^{-2E_*}$, $E_*=2(\tau+1)(\mu+1)+N+4$), so the dependence on the iteration index is nil; the constant $c'$ carries $2^{-O(\mu\tau)}$. The exponent is conservative in $L=a_s+1$ (which absorbs the logarithm in $T_0$) and in the maximum of three exponents.

\paragraph{The numerical check.} The scripts are in \texttt{.drafts/P/bp-scratch-pfix/} (\texttt{sched.py}, \texttt{rand.py}, \texttt{budget.py}; mpmath, 60 digits; the bump-function $L^1$ norms behind $c'_{N,\mu}$ are \emph{approximations} by a trapezoidal rule, not certified constants; the rest uses exact $\LamD\delta=\coth(\pi\delta)^N-1$, $d_\nu$, $E_\nu$, not the bounds of the proofs). \texttt{sched.py} builds the schedule of Definition~\ref{bpP:p11a:schedule} with $C_1\to\max\{1,C_1\}$, with $M$ and $M_0$ as separate inputs, $\gamma=\gamma_*/2$, $\varepsilon=\varepsilon_{\mathrm{adm}}$, and checks, for $\nu<90$: (C1)--(C3); $\sum d_\nu\le s/4$, $r_\nu\ge s/2$; $\lVert D\omega_\nu\rVert\le2C_1$; the recursion; the finite-stage tails $S_m\le\gamma T_m^{-\tau-1}$ for \emph{all} $m$, the infinite parts of the sums being bounded by the exact geometric envelope $\ell_{n-1}\varrho/(1-\varrho)$; the closeness sum (envelope ratio $g_3\le2^{-5}$); $Q_0\le\varpi/2$, $Q_0\Xi\le2^{-10}$; the budgets (E2) \emph{against the stated $2^{-12}$} (both sums), and (E3) against $2^{-10}$ (tail ratios checked $\le2^{-4}$). Four cases $(N,\mu,\rho,\tau;C_1,\vartheta,s,\beta,M{=}M_0)=(1,2,1,2;2,.5,.5,.3,3)$, $(2,3,2,3;2,.3,.3,.2,4)$, $(4,4,2,8;1,.2,.2,.1,2)$, $(6,5,2,15;1,.1,.1,.05,2)$: all inequalities hold; $\varepsilon_{\mathrm{adm}}\approx2.4\cdot10^{-86},\,1.1\cdot10^{-158},\,4.2\cdot10^{-442},\,1.9\cdot10^{-1036}$, with binding branch (Remark~\ref{bpP:p11a:adm-branches}) the quadratic one in the first case and the closeness one in the other three. \texttt{rand.py}: $40$ random data sets with independent $M,M_0$, $C_1$ below and above $1$: no failure; the quadratic branch binds in $17$ and the closeness branch in $23$ runs, the first never. \texttt{budget.py}: along families with a single rate equal to $1$ and along the family with all eight rates equal to $1$, for $(N,\mu,\rho,\tau)=(1,2,1,2)$ and $(2,3,2,3)$, the measured log--log slopes of $\varepsilon_{\mathrm{adm}}$ never exceed the exponent $B$ of Proposition~\ref{bpP:p11a:adm-budget}; they equal $B$ for $\beta,M,C_1,\varepsilon^*$ (and approach it for $\vartheta,d$), e.g.\ $14$ for $\beta$ in both cases ($=B_0$ of Corollary~\ref{bpP:p11a:adm-budget-beta}), $8$ and $9$ for $M$; for $s$ the measured slope ($12.2$, $24.6$) is below $B$ ($24$, $48$), reflecting the conservative $L=a_s+1$; for all rates $1$: $63.2\le77$ and $92.6\le116$. These are tests of the formulas; the proofs are Proposition~\ref{bpP:p11a:adm-choice}, Lemma~\ref{bpP:p11a:exported-sums} and Proposition~\ref{bpP:p11a:adm-budget}.

%%%%%%%%%%%%%%%%%%%%%%%%%%%%%%%%%%%%%%%%%%%%%%%%%%%%%%%%%%%%%%%%%%%%%%%%%%%%%%%%
\subsection{Sources, deviations, errata, what is reused}\label{bpP:p11a:notes3}

\paragraph{Sources read.} Rus01 \S13.1--13.4 (pp.~156--158; the printed formulas (13.1.2), (13.1) are partly garbled in the text layer of the pinned PDF: only $\kappa=24/25$, $\tau_0=1/9$ (read from the rendered page; the text layer drops the denominator) and the structure $\varepsilon_\nu=\varepsilon_0\delta^{\kappa\nu}$, $t_\nu=t_0\delta^{\tau\nu}$, $T_\nu=\Psi^{-1}(\Psi(T_0)\delta^{\tau\nu})$ are used), \S14.1--14.8 (pp.~165--170), \S16.1--16.7 (pp.~178--179), \S20.5--20.7 (pp.~201--203). Not read: \S\S13.5--13.12 (one-step remainder; recomputed in File~11a-2), \S15 (convergence; File~11a-4), \S\S19--20.4 (smoothing, chains; replaced). Files 11a-1, 11a-2; \texttt{reviews/p11a2-review-1.md} (P11A2-01, -03, -04); \texttt{PLAN-11a.md} (deviation table, R1, R2, R5).

\paragraph{Deviations (D-numbers as in PLAN-11a.md; D-16 to D-18 new).}
\begin{itemize}
\item \textbf{D-3} (power divisor): $D_\nu=\gamma T_\nu^{-\tau}$ (not $\gamma\Phi(T_\nu)$ with $\Psi=T^{-d}\Phi$); no coupling of $T_\nu$ to $\varepsilon_\nu$ through $\Psi^{-1}$, so Rus01's exponent bookkeeping $\kappa-\sigma-\tau_0>\ldots$ and the integral condition 13.1.6 do not occur (Lemma~\ref{bpP:p11a:strip-budget}).
\item \textbf{D-7} (chain extension): mollified cutoffs; verified, Remark~\ref{bpP:p11a:d7-verdict}. Condition $g_3\le2^{-5}$ instead of $\tau\mu<\tau_0$.
\item \textbf{D-9} (explicit $\varepsilon_{\mathrm{adm}}$): Definition~\ref{bpP:p11a:adm-params}.
\item \textbf{D-15} (strip bookkeeping): $\sum(d_\nu+\delta_\nu+2\sigma_\nu)\le s/2$, condition (A6).
\item \textbf{D-16} (decay ratio and growth ratios): $\varpi=2^{-2E_*}$ instead of $\delta^\kappa$ with $\kappa=\frac{24}{25}$; $\Theta=2$, $\zeta=\frac12$ free; the price is the constants $2^{-O(\mu\tau)}$ in $\varepsilon_{\mathrm{adm}}$.
\item \textbf{D-17} (collar): $t_\nu=t_02^{-(\tau+1)\nu}$ with $t_0\le\gamma T_0^{-\tau-1}/(4C_1)$, so that $t_{\nu+1}T_\nu\lVert D\omega_\nu\rVert\le D_\nu/2$ (Rus01 (13.4): $t_0\le\gamma\Psi(T_0)\vartheta/(8(2p^2+p)(3C_\bullet+3)^{2p^2+p})$, p.~158).
\item \textbf{D-18} (mode-wise test): $K_{\nu+1}$ tests the shell $\Sh_\nu$ with margin $2\gamma\lvert k\rvert_1^{-\tau}$ against $\omega_\nu$ itself (not the cutoff-wide $\frac12\gamma\Phi(T_\nu)$ of Rus01 (16.1)); the finite-stage invariant (Lemma~\ref{bpP:p11a:finite-divisor}) gives the cutoff-wide divisor.
\end{itemize}

\paragraph{Suspected errata.} None new in the pages read. (The text layer of Rus01 p.~156 is garbled at (13.1.2) and (13.1); not a defect of the source.) Defect inherited from File~11a-1: the endpoint and the limit-based form of \ref{bpP:p11a:I3} (Remark~\ref{bpP:p11a:i3-corrected}), already noted by the review of File~11a-2.

\paragraph{Doubts about frozen statements.} None: no frozen declaration is touched. Note for the consumers: the geometric scheme gives $\varepsilon_{\mathrm{adm}}$ polynomial but with a large degree; Part~11's finite preparation order must exceed the exponent $B_0$ of Corollary~\ref{bpP:p11a:adm-budget-beta}; its interface item (K5) must list $M_0$, $d$ and $C_1\ge1$ among the fixed data.

\paragraph{External facts used.} \ref{bpP:p11a:y-log}--\ref{bpP:p11a:y-geom}; Files~11a-1, 11a-2's lists; $\Kref{bp:p1:uniform-limits}$, $\Kref{bp:p1:exp-sums}$, $\Kref{bp:p1:decay}$.

\paragraph{Forward references.} Part~11a-4 (limit, tori, displacement, family): \ref{bpP:p11a:O3}, \ref{bpP:p11a:O4}, Lemma~\ref{bpP:p11a:exported-sums}; Part~11a-5 (budget): Definition~\ref{bpP:p11a:adm-params}, Theorem~\ref{bpP:p11a:parameter-selection}, Proposition~\ref{bpP:p11a:adm-budget}, Corollary~\ref{bpP:p11a:adm-budget-beta}. No unwritten label is cited except those.

\paragraph{Label index.} \texttt{sel-data}, \texttt{geom-recursion}, \texttt{schedule-consts}, \texttt{schedule-ratios}, \texttt{schedule}, \texttt{strip-budget}, \texttt{schedule-props}, \texttt{smooth-cutoff}, \texttt{finite-extension}, \texttt{d7-verdict}, \texttt{initial}, \texttt{adm-params}, \texttt{adm-choice}, \texttt{finite-divisor}, \texttt{i3-corrected}, \texttt{finite-measure}, \texttt{inductive-lemma}, \texttt{parameter-selection}, \texttt{exported-sums}, \texttt{adm-branches}, \texttt{adm-budget}, \texttt{adm-budget-necessary}, \texttt{adm-budget-beta}.
\par\endgroup
